\documentclass[11pt,letterpaper]{article}
\usepackage[T1]{fontenc}
\usepackage[utf8]{inputenc}
\usepackage{lmodern,microtype}
\usepackage[margin=.92in,headheight=15pt]{geometry}
\usepackage{amsmath,amssymb,amsthm,mathtools}
\usepackage{booktabs,longtable,array,enumitem,xcolor,fancyhdr,listings,needspace}
\usepackage{xurl}
\usepackage[colorlinks=true,linkcolor=blue!40!black,citecolor=blue!40!black,urlcolor=blue!45!black]{hyperref}
\usepackage{bookmark}
\definecolor{ink}{RGB}{24,48,65}
\definecolor{paperblue}{RGB}{246,249,251}
\hypersetup{pdftitle={Constructing Computation Sources: A Smaller Complete iPRA Checker},pdfauthor={Research prepared for Vladimir Reshetnikov},pdfsubject={Intuitionistic primitive recursive arithmetic, proof checking, bi-interpretability}}
\setlength{\parskip}{.3em}
\setlength{\emergencystretch}{3em}
\setlist{itemsep=.2em,topsep=.4em,leftmargin=1.7em}
\newtheorem{theorem}{Theorem}[section]
\newtheorem{lemma}[theorem]{Lemma}
\newtheorem{proposition}[theorem]{Proposition}
\newtheorem{corollary}[theorem]{Corollary}
\theoremstyle{definition}\newtheorem{definition}[theorem]{Definition}
\theoremstyle{remark}\newtheorem{remark}[theorem]{Remark}
\newcommand{\ipra}{\mathsf{iPRA}}
\newcommand{\Kcal}{\mathcal K}
\newcommand{\Ccal}{\mathcal C}
\newcommand{\FV}{\operatorname{FV}}
\newcommand{\file}[1]{\expandafter\nolinkurl\expandafter{\detokenize{#1}}}
\lstset{basicstyle=\ttfamily\footnotesize,columns=fullflexible,keepspaces=true,
 breaklines=true,showstringspaces=false,frame=single,rulecolor=\color{black!15},
 backgroundcolor=\color{paperblue},xleftmargin=.3em,xrightmargin=.3em}
\pagestyle{fancy}\fancyhf{}
\fancyhead[L]{\small Constructing computation sources}
\fancyhead[R]{\small Intuitionistic primitive recursive arithmetic}
\fancyfoot[C]{\thepage}
\title{\color{ink}\textbf{Constructing Computation Sources}\\[.5em]
\large A smaller complete iPRA checker and a finite-range obstruction}
\author{Research prepared for Vladimir Reshetnikov}
\date{November 10, 2026}
\begin{document}
\maketitle
\begin{abstract}
We give an explicit proof presentation of full first-order intuitionistic
primitive recursive arithmetic, with a standalone checker strictly below
the observed complete-checker minima in the nine current ProveIt draft
archives. Under the unchanged whole-program normalizer, the selected
checker has \textbf{279 tokens, 314 syntax-tree nodes, and 682 bytes},
compared with the current coordinate envelope \textbf{281/315/685}.

The mathematical change is small and precise. A computation certificate
supplies a predecessor and a base/step flag; the checker constructs the
source counter. The predecessor image is checked only if the update uses
it. This replaces an eager validation obligation by a construction
invariant and a state-independence argument. The change preserves the
native theory and its whole-domain bi-interpretation with iPRA. An
explicitly attributed appendix reproduces the full inherited arithmetic
and logical bridge.

A separate new result settles a current-draft question negatively:
restricting paired-introduction witnesses to zero or variables and
closure-specialization targets to one or variables loses native theorems.
Adding any fixed finite set of numeral targets does not repair the loss.

These are a reproducible implementation reduction and a mathematical
obstruction theorem. They do not establish an absolute simplest
axiomatization or the significant breakthrough requested. The source-size
gain is modest; the arithmetic basis and the logical reconstruction are
inherited. The proofs are written mathematical arguments, not a
Lean or Rocq verification.
\end{abstract}
\clearpage
\setcounter{tocdepth}{1}\tableofcontents
\clearpage

\section{The precise improvement}
\label{sec:comparison}

The comparison is pinned to ProveIt commit
\begin{center}\small\texttt{4a04139595eadc3e8818f1df2a94838ed4851d83}.\end{center}
The \file{docs/drafts} tree is
\texttt{cec0cd82d69646386c802841068c17ff76d188fc}.
It contains nine direct ZIP archives\cite{snapshot}. Their exact bytes,
Git blob hashes, SHA-256 hashes, and member inventories are included in
the package. The source scan measured 261 extracted Python occurrences
with 164 distinct contents. An additional scan of ten nested provenance
archives measured 154 historical occurrences, with 109 distinct contents;
none improves the current complete-checker envelope.

Files smaller than the envelope were inspected before classification:
build, measurement, experiment, and verification scripts are not complete
theorem checkers. Complete historical checkers inside reference folders
do count. The comparison is with the present combined corpus, not the
older minima quoted by individual reports.

\begin{center}\small
\begin{tabular}{lrrr}
\toprule
Entire standalone program & Tokens & AST nodes & Bytes\\
\midrule
Generated assumptions, principal\cite{generated} &281&318&685\\
Generated assumptions, AST profile\cite{generated}&283&315&710\\
Current coordinate envelope&281&315&685\\
\midrule
Constructed computation, explicit assignments&277&317&684\\
Constructed computation, AST profile&279&314&709\\
\textbf{Constructed computation, selected profile}&\textbf{279}&\textbf{314}&\textbf{682}\\
\midrule
Old generator with the same initialization profile&283&315&683\\
\bottomrule
\end{tabular}
\end{center}

The selected single program strictly improves all three current minima:
two fewer tokens, one fewer AST node, and three fewer bytes. The percentage
reductions against those separate minima are approximately $0.71\%$,
$0.32\%$, and $0.44\%$. They are small. The three minima are attained
by different old programs, whereas the selected new program beats all
three simultaneously.

The arithmetic generator itself changes the explicit program from
$281/318/685$ to $277/317/684$. The selected profile additionally folds
initialization into an existing default expression and folds the input
binding into its first call. Those two Python idioms are inherited
source-layout techniques, not new mathematical principles. Applying the
identical profile to the old generator yields $283/315/683$: the genuine
generator comparison remains four tokens, one AST node, and one byte.
The table separates these effects so that the improvement cannot be
misattributed to a new arithmetic foundation.

\subsection{What the metric includes}

The unchanged normalizer removes genuine docstrings, parses the complete
module, applies \texttt{ast.unparse} followed by a newline, and counts
significant tokens, all AST nodes, and UTF-8 bytes. Measurements use
CPython 3.12.14. Initial state, JSON input, traversal, scope checks,
inference dispatch, and the complete ACCEPT/REJECT driver are included.
Every counted checker imports only the standard JSON library. Producers,
converters, and the numerical testing decoder are not trusted by it.

The Python runtime and standard library are the shared platform, as in
the drafts. No interpreter, axiom table, theorem-specific library, or
unchecked executable certificate is moved outside the counted file.
The selected and old programs have the same outer JSON container,
four-field packet arity, and verdict convention. Computation packets use
the new predecessor/mode semantics and require translation.
In particular, rejected inputs print \texttt{REJECT}, exit nonzero, and
may produce a traceback; accepted inputs print \texttt{ACCEPT} and exit
zero. The program does not use assertions whose behavior changes under
Python optimization.

The observed sizes are upper bounds on an implementation optimum. They
are not mathematical lower bounds for all theories or all programming
languages. This report meets the operational comparison with the supplied
complete programs; it does not prove global minimality or claim that a
one-byte algorithmic reduction constitutes a major simplification of
arithmetic.

\section{The complete native axiomatization}
\label{sec:native}

The target is \emph{full first-order} iPRA: intuitionistic logic with actual
numerical equality, genuine falsity, function symbols for fixed finite
primitive-recursive descriptions, their defining equations, the usual
successor boundary laws, and quantifier-free induction. Atomic-induction
presentations of iPRA also occur in the literature\cite{salehi}; here the
reduction from quantifier-free induction is proved internally in the
reproduced bridge. No equivalence with a merely quantifier-free theorem
set is substituted for bi-interpretability.

\subsection{Numerical and formula syntax}

There is one numerical sort. The numerical terms are
\[
 t::=x\mid 0\mid St\mid I(n,a;v.h).
\]
Only $h$ binds $v$; $n,a$ do not. Other parameters of $h$ remain fixed.
Each body is a fixed finite earlier description. Equivalently, at each
external syntactic stage introduce a first-order symbol $i_h(n,a,\vec p)$
for each earlier term $h(v,\vec p)$. There is no function quantifier and
no uniformly total evaluator for all primitive-recursive descriptions.

The two computation patterns are
\begin{align}
 I(0,a;v.h)&=a,\tag{I0}\label{eq:I0}\\
 I(Sn,a;v.h)&=h[I(n,a;v.h)/v].\tag{IS}\label{eq:IS}
\end{align}
The formula grammar is
\[
 A::=(s=t)\mid Q_yA\mid D_{x,y}A,
\]
with distinct paired bound variables. A context is a finite reusable
list of these formulas. Substitution is capture-free; binders are renamed
before substituting. The designated absurdity is $\bot_*=(0=S0)$.
The ordinary readings, whose adequacy will be proved rather than assumed,
are
\begin{align*}
 (Q_yB)^\circ&=B^\circ(0)\to\forall y\,B^\circ(y),\\
 (D_{x,y}B)^\circ&=\exists x\bigl(B^\circ(x,0)\to
                                  \forall y\,B^\circ(x,y)\bigr).
\end{align*}

\subsection{Every primitive rule}

The presentation $\Kcal$ has the following complete inventory.
Hypothesis use, equality transport, and closure use are
\[
 \frac{}{\Gamma\vdash A}\;(A\in\Gamma),\qquad
 \frac{\Gamma\vdash s=t\quad\Gamma\vdash F(s)}
      {\Gamma\vdash F(t)},\qquad
 \frac{\Gamma\vdash Q_yB(y)\quad\Gamma\vdash B(0)}
      {\Gamma\vdash B(t)}.
 \tag{Hyp, Tr, QE}
\]
The motive $F$ is any fixed native formula template. Transport does not
license induction on $F$.

Paired elimination is
\[
 \frac{\Gamma\vdash D_{x,y}B(x,y)\qquad
       \Gamma,Q_yB(a,y)\vdash A}
      {\Gamma\vdash A}.
 \tag{DE}
\]
Here $a$ is absent from $\Gamma,A$ and the fixed parameters of $B$.
Paired introduction has two forms:
\[
 \frac{\Gamma,B(t,0)\vdash B(t,Sz)}
      {\Gamma\vdash D_{x,y}B(x,y)},
 \qquad
 \frac{\Gamma,E(t,z)\vdash E(t,Sz)}
      {\Gamma\vdash D_{x,y}E(x,y)}.
 \tag{DI-case, DI-progress}
\]
The second matrix $E$ must be a \emph{literal equation}. A formula with a
$Q$ or $D$ at its root is not literal. In both forms $z$ is absent from
$\Gamma,t$ and all fixed matrix parameters. Displayed local hypotheses
are discharged. Case introduction is allowed for every native matrix.

Finally, numerical computation is supplied as a temporary assumption:
\[
 \frac{\Gamma,E\vdash A}{\Gamma\vdash A}
 \quad\left(E\text{ is an instance of I0 or IS}\right).
 \tag{Comp}
\]
This is two finitely specified computation schemas, whose instances are
arbitrary finite terms with the stated bindings. There are no additional
logical or nonlogical axioms, primitive reflexivity rules, coverage
axioms, separation axioms, unrestricted induction rules, or hidden
first-order proof engines.

The reference calculus $\Ccal$ instead has I0 and IS as theorem leaves
and the same five logical rule families. Both presentations prove exactly
the same native sequents: Comp is cut with a computation theorem, and a
computation theorem is obtained in $\Kcal$ by Comp followed by Hyp.
Weakening, exchange, contraction, renaming, numerical substitution, and
cut are admissible by induction on finite derivations. Rename local
parameters before substitution or cut; substitution preserves the root
constructor and hence the literal-equation induction restriction.

\section{The new computation-certificate principle}
\label{sec:principle}

The current leading draft supplies an alleged iterator source
$[c,a,h]$. In a successor step it expands $c$ and eagerly checks the
predecessor iterator before substituting it into $h$. This guards against
a non-successor counter that happens to have a truthy serialization.
The old check is important: simply deleting it admits a false equation.

The new certificate supplies the \emph{predecessor} $p$, the initializer
$a$, the update $h$, and an explicit mode $e$. The checker itself forms
the source:
\[
\begin{array}{c|c}
e=0 & I(0,a;v.h)=a\\[2pt]
e=1 & I(Sp,a;v.h)=h[I(p,a;v.h)/v].
\end{array}
\]
The mode occupies an already existing, previously unused computation
field. It does not lengthen the packet. The source has a genuine
successor constructor because the checker constructs it. In particular,
the successor-of-zero case uses $p=0,e=1$; there is no restriction to
nonzero predecessors or any special zero-alias convention.

The two lines implementing the new generator are
\begin{lstlisting}[language=Python]
n, t, r = s
h = [[u and [n], t, r], i(r, s, 1) if u else t]
\end{lstlisting}
Here $s=[p,a,h]$ and the packet's $u$ is the mode. The local $t,r$
become the initializer and update. Only canonical flags 0 and 1 are
needed for completeness; other raw values are handled as described below.

\subsection{Why validation can be delayed}

The opener $i$ checks a substituted numerical image whenever an exposed
numerical coordinate actually uses that image. Therefore the old eager
validation can be separated into two distinct obligations:
\begin{enumerate}
\item the source counter is a successor;
\item a recursive image used in the result is a valid numerical term.
\end{enumerate}
Construction establishes the first. Demand-driven substitution establishes
the second. If the update ignores its state, the recursive image is absent
from the generated equation and needs no separate validation.

There is a further scope issue: the common continuation introduces a new
parameter $w$, and a raw generated equation can depend on it. Its zero
instance, rather than the equation at every $w$, is what must be proved.

\begin{lemma}[Known-instance discharge]
\label{lem:known-instance}
If $w\notin\FV(\Gamma,A)$, $\Gamma\vdash H(0)$, and
$\Gamma,H(w)\vdash A$, then $\Gamma\vdash A$.
\end{lemma}
\begin{proof}
Rename eigenparameters in the continuation to avoid capture. Substitute
zero for $w$, leaving the retained context and goal unchanged. This gives
$\Gamma,H(0)\vdash A$. Cut with the given proof of $H(0)$.
\end{proof}

\begin{theorem}[Constructed computation supply]
\label{thm:new-supply}
Every equation $H(w)$ successfully generated by the new arithmetic
branch has a zero instance provable from I0 or IS, with equality
substitution. Consequently its common continuation is admissible in
$\Ccal$ and hence in $\Kcal$.
\end{theorem}
\begin{proof}
In base mode, after setting $w=0$, the source initializer and its copy on
the right agree by structural zero completion. I0 gives the equation.
The predecessor payload is unused and has no object-language meaning
in this branch.

In step mode, first suppose that opening the update requests its recursive
image. It then validates $[p,a,h]$ at numerical root. Successful numerical
validation is idempotent, so doing it on demand gives the same substituted
value as validating eagerly. The extra dummy zero slot beneath the
constructed successor cannot introduce a meaningful new dependency:
numerical-root validation of the predecessor image would reject that
dependency. The two occurrences of the predecessor therefore agree, and
the zero instance is IS.

If the update never requests its image, it is independent of the outer
state on every meaningful numerical path. Genuine internal variables
are protected; a meaningful outer-state reference would request the
image. Numerical binary annotations are semantically zero and contribute
no hidden state dependence. The common opening still checks the entire
source iterator. After zero completion it has the form
$I(Sp_0,a_0;v.h_0)=h_0$, where $h_0$ is independent of $v$. IS proves
this equation without needing an unused recursive image.

The old goal and context were checked before allocating $w$, so they
have no meaningful occurrence of it. Lemma~\ref{lem:known-instance}
discharges the actual generated assumption. Appendix
\ref{sec:constructed-supply-audit} gives the full raw-coordinate argument.
\end{proof}

\begin{corollary}[No change to the native theory]
The new certificate calculus proves exactly the closed native theorems
of $\Kcal$. It preserves the same sequents and the same interpretation
composites, rather than merely preserving standard numerical truth.
\end{corollary}
\begin{proof}
Soundness follows from Theorem~\ref{thm:new-supply} and the unchanged
logical-rule checks. For completeness, encode each Comp instance with
its mode and predecessor; supply the generated equation and close its
use by Hyp. Other instructions and their eigenparameter allocation are
unchanged. Canonical numerical syntax uses rigid labels for passive
parameters and zeros in dummy positions. Both computation equations,
including the step with zero predecessor, are generated exactly.
\end{proof}

\section{Recovering logic and arithmetic}
\label{sec:bridge-summary}

This section summarizes the inherited construction; the complete
arguments appear in Appendices~\ref{sec:choice-introduction}--\ref{sec:ordinary-bridge}.
Their provenance and unchanged source hashes are recorded. The new
contribution is the conservative certificate generator above, not the
following arithmetic bootstrap.
The inherited logic appendix presents $D$ formulas with auxiliary
clauses. The present nested-$Q$ syntax is covered because $Q_yB$ is
interderivable with $D_{v,y}B$ for unused $v$: case DI gives one direction,
and DE followed by Hyp gives the other. External induction eliminates
nested $Q$ occurrences, using the same uniform matrix-congruence
argument. Thus no extra logical connective is left without a round trip.

Put $C(n,a,b)=I(n,a;v.b)$. Then $C(0,a,b)=a$ and $C(Sn,a,b)=b$.
From computation and transport one derives reflexivity, symmetry,
transitivity, and congruence. DI-progress, DE, and QE give contextual
literal-equation induction. DI-case, DE, and QE give numerical cases
for arbitrary native formulas. These constructions retain arbitrary
assumptions, with the displayed eigenparameter conditions.

For a formula $A$, form a gate $G_zA$ recursively. At an equation leaf
replace $s=t$ by $C(z,0,s)=C(z,0,t)$, and continue beneath binders.
Then $G_0A$ is provable and $G_{Sz}A$ is interderivable with $A$.
Padding by a vacuous paired binder allows two formulas of different
depths to be aligned. A recursively constructed selector $M_z(A,B)$
then satisfies
\[
 M_0(A,B)\dashv\vdash A,\qquad M_{Sz}(A,B)\dashv\vdash B.
\]
The scalar parity term $p(z)=I(z,0;v.C(v,1,0))$ has $p(1)=1,p(2)=0$.
With fresh bound variables, the following macros recover the full
intuitionistic interface:
\begin{align*}
 A\Rightarrow_*B&=D_{u,z}M_z(A,B),\\
 A\wedge_*B&=D_{u,z}G_zM_{p(z)}(A,B),\\
 \exists_*x\,A(x)&=D_{x,y}G_yA(x),\\
 A\vee_*B&=\exists_*x\,M_x(A,B),\\
 \forall_*x\,A(x)&=A(0)\wedge_*D_{u,z}G_zA(z),\\
 \bot_*&=(0=S0).
\end{align*}
The coordinate $u$ is unused. The selector gives implication rules;
the gate gives existential rules; numerical cases and selectors give
disjunction; the two parity values give conjunction projections.
Coverage gives universal elimination in arbitrary retained contexts.
Under $0=S0$, the selector and transport give every equation; structural
formula recursion then gives explosion. No excluded-middle rule for
quantified formulas is introduced.

The ordinary realization $T_I$ has standard intuitionistic equality
logic, I0 and IS, contextual literal-equation induction, and the explicit
boundary sentences
\[
 \forall x(x=0\vee\exists y\,x=Sy),\qquad (0=S0)\to\bot.
\]
The falsity bridge is necessary: positive numerical equations alone
admit a one-element structure.

The arithmetic reconstruction has a fixed noncircular order. First an
internal scalar predecessor is built using numerical cases and equation
induction. Then truncated subtraction and a characteristic of actual
equality are constructed and proved correct for all source elements.
Each fixed quantifier-free formula has a characteristic equation, so its
induction follows. Finally pairing and projections are constructed and
general primitive recursion is compiled. Pairing is not assumed in the
predecessor argument.

If $\pi,L,H$ are the proved pairing and projections, a general recursion
$f(0,\vec p)=g(\vec p)$, $f(Sn,\vec p)=h(n,f(n,\vec p),\vec p)$ is
represented by
\[
\begin{split}
 V(n,\vec p)&=I\bigl(n,\pi(0,g^*(\vec p));
 v.\pi(SL(v),h^*(L(v),H(v),\vec p))\bigr),\\
 f^*(n,\vec p)&=H(V(n,\vec p)).
\end{split}
\]
Equation induction proves $L(V(n,\vec p))=n$, and the projection
identities prove both defining equations. This is an external compiler
for each finite description, not an internal uniformly total evaluator.

\begin{theorem}[Whole-domain bi-interpretability and equiconsistency]
\label{thm:biinterp}
The ordinary realization of $\Kcal$ is bi-interpretable and equiconsistent
with full first-order iPRA. In each numerical composite, the definable
isomorphism is $x=y$ on the entire domain.
\end{theorem}
\begin{proof}
Interpret scalar iteration in iPRA by primitive recursion ignoring the
counter in its update; interpret each iPRA description in $T_I$ by the
scalar compiler. The preceding internal arithmetic proves all source
axioms and quantifier-free induction instances. The logical macros and
ordinary readings preserve all formulas and proof rules.

For either numerical composite, the double translation satisfies the
original defining equations. Literal-equation induction proves uniqueness
with all passive parameters fixed. External induction on a fixed finite
description therefore identifies the operation with its double
translation at every internal argument. The identity graph is internally
total, bijective, and operation-preserving. No quotient, definable cut,
or restriction to standard numerals is used.

For formula composites, the decisive native equivalences are
\[
 Q_yB\dashv\vdash B(0)\Rightarrow_*\forall_*yB(y),
\]
\[
 D_{x,y}B\dashv\vdash
 \exists_*x\bigl(B(x,0)\Rightarrow_*\forall_*yB(x,y)\bigr).
\]
DE proves the forward direction by exposing a witness and its closure;
the reverse uses existential elimination and case DI, with its successor
premise supplied by implication and universal elimination. Structural
induction gives all formula round trips. The explicit falsity bridge
preserves contradiction. Theorem~\ref{thm:new-supply} transfers the
interpretations unchanged to the new certificate presentation. Finite
proof translations yield equiconsistency, with no claim that either
theory proves its own consistency.
\end{proof}

\Needspace{24\baselineskip}
\section{The entire selected checker}
\label{sec:checker}

\begin{lstlisting}[language=Python]
from json import load
i = lambda a, t=None, d=(k := 0): [i(b, t, d + b' \x02\x01\x04'[len(a)]) for b in a] if isinstance(a, list) else a if k <= a < d else d % 2 and i(t, d=1)

def c(a, g):
    global k
    if i(a) in g:
        return
    m, s, u, t = q()
    if m:
        c(m, g)
        if m in ([s, u], not s and [t]) and i(t, u) == a:
            a = i(t, s)
        h = m[2:]
    elif t:
        h = a = i(a[2], t, 2)
    else:
        n, t, r = s
        h = [[u and [n], t, r], i(r, s, 1) if u else t]
    k -= 1
    c(i(a, [k]), g + [i(h, k if m or len(a) % 2 == m else 0)])
try:
    c((q := load(open(0)).pop)(), [])
except:
    print('REJECT')
    raise
print('ACCEPT')
\end{lstlisting}

The assignment to $k$ in the lambda's default is evaluated once when the
function is created at module scope; it does not reset freshness during
calls. Likewise the input assignment initializes $q$ before its first
call. The explicit-assignment implementation has the same successful
traces and is supplied for easier inspection. The descriptive reference
implementation uses meaningful names and separates the branches. None
of these files calls a proof producer or a theorem library.

\subsection{Input, tree decoding, and scopes}

Input is a JSON array. Its last item is the requested closed formula;
the preceding four-field packets $[m,s,u,t]$ are popped from the end.
Hypothesis closure consumes no packet. An unread prefix is inert, as in
the comparison programs. A malformed, exhausted, or unsuccessful input
rejects. Physical Python recursion or memory limits may reject large
valid certificates; mathematical completeness concerns unbounded
resources on finite canonical certificates.

\begin{center}\small
\begin{tabular}{cllc}
\toprule
Array length & Formula reading & Numerical reading & Depth added\\
\midrule
0 & truth $0=0$ & zero & no children\\
1 & $Q$ & successor &2\\
2 & equation & zero with checked annotations &1\\
3 & $D$ & scalar iterator &4\\
\bottomrule
\end{tabular}
\end{center}
Larger arrays reject when traversed. Let $\rho$ be the environment of
protected numerical slots. The main decoder clauses are
\begin{align*}
 T_\rho([p])&=S T_{0,\rho}(p),&
 F_\rho([B])&=Q_yF_{y,\rho}(B),\\
 T_\rho([p,q])&=0,&
 F_\rho([s,t])&=(T_\rho(s)=T_\rho(t)),\\
 T_\rho([n,a,h])&=I(T_{0,0,\rho}(n),T_{0,0,\rho}(a);
                         v.T_{v,v,\rho}(h)),\\
 F_\rho([P,U,B])&=D_{x,y}F_{y,x,\rho}(B).
\end{align*}
The two padding fields of $D$ and the annotations of a numerical binary
node are checked, although they contribute no mathematical data. In
particular the scalar iterator has two aliases for its one state, not
an extra recursion-counter input.

At formula depth $2q$ and numerical depth $2q+1$ there are $q$ protected
slots. Positive coordinates $2j-1,2j$ name the same $j$th slot, nearest
first. Scalar zero denotes numerical zero; negative labels name already
allocated rigid parameters. The floor $k$ starts at zero and decreases
at each continuation. The opener preserves labels in $[k,d)$, inserts a
checked numerical image at an exposed numerical coordinate, and maps an
exposed formula leaf to truth. Missing meaningful images reject.
Canonical inputs use integers and arrays. Accepted finite fractional
aliases have the inherited floor-coordinate meaning; noncanonical data
do not add new numerical operations. Scope safety and the common-motive
opening lemma are structural facts about this traversal.

\subsection{All command branches}

For nonempty $m$, the checker first proves the major $m$. Matched equality
or closure endpoints and the target equality implement Tr or QE. They
change the goal to the appropriate minor and add only truth to the
context. An inapplicable Eq/Q request leaves the goal unchanged and adds
only truth. A three-field $D$ major adds its opened $Q$ closure, giving
DE. An unsuccessful major never falls through as if it had succeeded.

For false $m$ and true $t$, the command is DI with witness $t$.
Canonical marker $m=[]$ selects the zero-case assumption; marker $m=0$
selects the current hypothesis only when the matrix is a literal
equation, and otherwise selects the zero case. A fresh successor is the
new goal. Thus arbitrary compound induction is not silently introduced.

For false $m$ and false $t$, the new arithmetic generator uses
$s=[p,a,h]$ and mode $u$. It constructs the source and supplies the
generated equation. Theorems above justify all successful raw cases,
including those whose generated equation has only a known zero witness.

Finally every non-hypothesis command decreases the floor, opens the new
goal and hypothesis, and takes one common continuation. The prior context
and goal cannot contain the fresh label meaningfully. Every successfully
added context formula is an array. Hence a scalar goal cannot terminate
successfully by exploiting the truth placeholder: it either rejects or
continues consuming packets. These invariants ensure that every successful
major has an intended formula shape.

\section{A new obstruction to simultaneous restrictions}
\label{sec:obstruction-summary}

The named-specialization draft simulates arbitrary QE targets using
equality naming and separately discusses restricted introduction
witnesses. Whether both restrictions can be imposed together is a
different question\cite{named}. It has a negative answer for the native
rules with no extra arbitrary-substitution rule.

\begin{theorem}[Simultaneous restriction loses theorems]
Restrict DI witnesses to zero or in-scope variables, and QE targets to
one or in-scope variables. Then
\[
                 D_{x,y}\bigl(C(y,1,x)=1\bigr)
\]
is unprovable, although it is a theorem of the unrestricted calculus.
Adding any fixed finite set of numeral QE targets still leaves some
unrestricted native theorems unprovable.
\end{theorem}
\begin{proof}[Proof of the first assertion]
Evaluate numerical terms in the standard natural numbers, give every
free variable value zero, and define a proof invariant by
\[
 [Q_yB]=[B(0)]\Rightarrow[B(1)],\qquad
 [D_{x,y}B]=[B(0,0)]\Rightarrow[B(0,1)].
\]
An allowed DI witness has value zero; its fresh step parameter has value
zero, so its successor has value one. Both DI modes preserve the
displayed implication. DE's new closure has the same truth value as
its major. Variable-target QE supplies the base, and target-one QE is
implication elimination. Equality transport and computation preserve
the invariant; known-zero-instance Supply preserves it by choosing its
fresh parameter to be zero.

The displayed formula evaluates to true implying false: its zero
instance is $C(0,1,0)=1$, while its one instance is $C(1,1,0)=1$.
In the unrestricted calculus, case DI with witness 1 proves the formula
immediately from $C(Sz,1,1)=1$. Appendix
\ref{sec:finite-restriction} gives the finite-range extension.
\end{proof}

This establishes failure of identity-language theorem equivalence for
the specific simultaneous restrictions. It does not rule out all other
interpretations, all restricted calculi, or all alternative domains. In
particular an additional arbitrary-term substitution rule would change
the system and invalidate the claimed restriction.

\section{Verification, reproduction, and limits}
\label{sec:evidence}

% BEGIN evidence/verification_summary.tex
The frozen compact and descriptive checkers passed all 144 complete-program certificate and malformed-input runs, with 3 successful targeted mutation controls. The 33 certificate cases were run on both programs under ordinary Python and \texttt{-O}; the remaining runs test malformed input. All 28 legacy cases are accounted for: 26 canonical translations, 24 already identical to replayed packets, and four additional successful CLI runs for the other two. The two obsolete successor-wrapper negatives were replaced explicitly rather than mislabeled under the new packet meaning.

The independent generator audit attempted 76,960 packets, obtained 57,640 successful generated equations, and checked 518,760 standard-model valuations at the fresh zero witness. It found no unequal endpoints and had no resource-bound skips. The compact source SHA-256 is
\begin{center}\footnotesize\file{bf17dcc00ea37d8e2e251c86b3ab53adca55ed403ca445bc22b9ee1564d5b540}.\end{center}


% END include


The tests execute the complete standalone programs, including their
input and verdict drivers, under ordinary and optimized Python. The
examples exercise computation with used and unused state, transport,
closure specialization, nonvacuous witnesses, compound case introduction,
and contextual equation induction. Negative controls exercise failed
majors, forbidden compound progress, future labels, malformed sources,
and the false-equation exploit motivating the old eager guard. New
constructor packets are checked against their new meaning, not assigned
the semantics of an old packet with the same raw fields.

An independent finite numerical decoder evaluates zero instances of
generated assumptions, including dummy slots and preexisting rigid
parameters. These are falsification tests and reproducibility evidence.
They do not prove the universal soundness, completeness, or
bi-interpretation theorems; the written structural proofs establish
those claims. The metatheory and Python runtime are not machine-verified.

The package contains the selected checker, two explicit profiles, a
descriptive checker, an untrusted proof producer and converter,
certificates, the unchanged source-measurement utility, source archives,
hash manifests, test receipts, the complete LaTeX source, and this PDF.
Run \texttt{python3 verify.py} for the finite reproduction and
\texttt{python3 build.py} to rebuild the PDF and standalone TeX file.
The original archives remain unchanged.

The construction retains the substantial inherited scalar-predecessor
and equality-reflection arguments. It establishes a small strict
implementation reduction, not a substantially simpler arithmetic
foundation. The separate obstruction identifies a specific incorrect
way to reduce the logical rules. Neither the source audit nor these
proofs establish global minimality.

\appendix
% BEGIN appendices/constructed_supply.tex
% Independent proof audit; includable in the final report.
\section{Computation sources certified by construction}
\label{sec:constructed-supply-audit}

The computation packet has a mode field $e$ and a triple $[p,a,r]$.
The old packet interface already had an unused field in its computation
branch; that field now carries $e$.  The kernel generates the source
counter itself.  The relevant code is
\begin{verbatim}
    n, t, r = s
    h = [[u and [n], t, r], i(r, s, 1) if u else t]
\end{verbatim}
Here \texttt{u} is the mode $e$, \texttt{s} is $[p,a,r]$, and
\texttt{t} is reassigned to the initializer $a$.  Canonical mode values
are $0$ for the base equation and $1$ for the successor equation.

For base mode the generated equation is
\[
                 I(0,a;v.r)=a.
\]
The submitted predecessor payload $p$ is unused in this case.  For step
mode the generated equation is
\[
                 I(Sp,a;v.r)=r[I(p,a;v.r)/v].
\]
The step counter is a singleton array because the kernel constructs
\texttt{[n]}.  It is never an untrusted purported successor.  Thus no
counter-shape validator is needed in this branch.

The expression \texttt{u and [n]} can return a noncanonical falsy mode
unchanged.  This causes no additional arithmetic case.  If the subsequent
common opening succeeds, a falsy numerical scalar or empty array denotes
zero.  Unsupported falsy JSON values reject when traversed.  A truthy
mode is used only as a branch tag and always produces the same singleton
constructor.  Its own JSON shape has no numerical meaning.

\subsection{Demanded images}

The opener validates an inserted image at numerical root whenever a
meaningful exposed numerical coordinate requests that image.  Consequently,
if opening $r$ demands the predecessor image $[p,a,r]$, the latter is
checked before insertion.  This validation is exactly the old eager
check, performed at the first point where its result can affect the
equation.  Numerical-root validation is idempotent on successful results:
structural recursion leaves permitted coordinates unchanged, rejects a
missing meaningful numerical image, and replaces only inessential
formula leaves in checked annotations.

There is one apparent binding mismatch to check.  In the constructed
source, the displayed successor contributes an additional dummy zero
slot to the environment for $p$.  The predecessor image has only the
iterator's usual two dummy slots.  Root validation of the predecessor
image excludes every meaningful reference to the extra outer slot.
All genuine internal slots and existing rigid names therefore have the
same meanings in both occurrences of $p$.  The update's two aliased
state slots are replaced by the validated predecessor image.  After
decoding, this gives precisely the scalar successor equation.

\subsection{Unused images}

If opening $r$ never requests its image, $r$ is independent of its outer
state on every meaningful numerical path.  Indeed, at opening depth
$d$, references to internal binders are below $d$, while references to
either removed state slot would request the image.  Paths beneath a
numerical binary annotation have no numerical meaning and do not affect
this argument.  Existing rigid names remain parameters.

The common continuation still opens the complete generated equation,
and in particular traverses its source iterator.  It therefore validates
the submitted predecessor, initializer, and update in the environments
where they actually occur.  A malformed unused predecessor image is
not itself a term used by the equation.  Let $w$ be the newly allocated
parameter.  After setting $w=0$ and completing exposed outer slots by
zero, write the source data as $p_0,a_0,r_0$.  The generated equation is
then
\[
 I(Sp_0,a_0;v.r_0)=r_0,
\]
with $r_0$ independent of $v$.  The scalar successor computation axiom
proves it, because substitution of the unused recursive value changes
nothing.  Neither a validation nor a dummy-slot comparison of that
unused raw recursive image is required.

\subsection{The base equation and the fresh parameter}

In base mode the source initializer occurs under two dummy zero slots,
whereas its copy on the right may expose those coordinates during the
common opening.  The standard zero-extension argument applies: substitute
zero for the new parameter, preserve genuine local binders and old rigid
names, and compare the two finite syntax trees structurally.  Their
decoded initializers agree.  The zero instance is therefore the base
computation equation.

If the previous name floor is $k_0$, the new name is $w=k_0-1$.
The original goal and retained assumptions were checked before this
allocation, so they have no meaningful occurrence of $w$.  Generated
source data may acquire $w$ at the later common opening.  The three
cases above establish a computation proof of the actual added equation's
zero instance, $H(0)$.  A successful continuation
\(\Gamma,H(w)\vdash A\) is discharged by renaming its later
eigenparameters, substituting $w=0$, and cutting with that computation
proof.  The old end sequent is unchanged.  The rule never asserts that
an arbitrary added equation $H(w)$ holds for every $w$.

\subsection{Canonical completeness without zero aliases}

Every ordinary base computation is represented by mode $0$, its actual
initializer and update, and an arbitrary unused predecessor payload
(canonically zero).  Every ordinary successor computation with
predecessor $p$ is represented by mode $1$ and the triple $[p,a,r]$.
In particular, a genuine zero predecessor uses \texttt{n=0} together
with \texttt{u=1}, and the generated source counter is exactly
\texttt{[0]}.  No truthy zero alias or additional equality proof is
needed.

As in the inherited canonical emitter, active passive parameters are
allocated rigid names, and dummy slots contain zero rather than
references to otherwise unused enclosing slots.  Thus inserting the
displayed successor introduces no new dependency, and canonical opened
updates have exactly their specified state substitutions.  The result
is the original computation equation in the original raw canonical
syntax.  Supplying it and immediately applying the hypothesis rule
simulates an ordinary computation leaf.  Other proof instructions and
their completeness transformations are unchanged, apart from the
already specified fresh-label accounting.

\paragraph{Precise scope.}
This is a new certificate format and a conservative computation generator.
An old packet whose source counter was malformed is not being newly
accepted with its old meaning.  The new packet supplies a predecessor,
and the checker constructs a genuine successor around it.  Its actual
generated equation is the one that must be used in the continuation.
For example, a binary numerical predecessor denotes zero and yields a
successor counter; it cannot trigger the false equation obtained by
mistaking that binary node itself for a successor.  The entire argument
is about finite syntax and ordinary computation axioms on the full
numerical domain.  It does not use numeral testing as a substitute for
the theorem, and does not claim an absolute smallest proof checker.

% END include

% BEGIN appendices/restriction_obstruction.tex
\section{A finite-range obstruction to simultaneous witness and specialization restriction}
\label{sec:finite-restriction}
Let $\mathcal C$ be the native Choice--DI calculus with fixed scalar terms
$0,S,I$, equations, binders $Q_yB$ and $D_{x,y}B$, and the rules specified in
\texttt{ipra\_named\_specialization}. Consider the simultaneous restriction
\[
\text{DI witnesses: }0\text{ or an in-scope variable},\qquad
\text{QE targets: }1\text{ or an in-scope variable}.
\tag{R}
\]
Neither arbitrary numerical substitution nor ordinary quantified logic
is silently added as a primitive rule.

\begin{theorem}
Restriction (R) strictly decreases the native theorem set. In particular,
with $C(y,p,q)=I(y,p;u.q)$ and $1=S0$, the formula
\[
D_{x,y}\bigl(C(y,1,x)=1\bigr)
\tag{N1}
\]
is a theorem of $\mathcal C$ but not of (R).
\end{theorem}
\begin{proof}
Evaluate numerical terms by their standard natural-number functions and
assign zero to every free numerical variable. Define native formula truth
recursively by ordinary equality and
\[
\begin{split}
[Q_yB]&=[B(y:=0)]\Rightarrow[B(y:=1)],\\
[D_{x,y}B]&=[B(x:=0,y:=0)]\Rightarrow[B(x:=0,y:=1)].
\end{split}
\]
These clauses specify a proof invariant, not an ordinary first-order
interpretation with unrestricted substitution. Formula substitution
respects equality of the numerical values of substituted terms, since
all binders are renamed away from those terms before substitution.
Thus equality transport and the computation schemas preserve validity.
Hypothesis use does so immediately.

An allowed DI witness has numerical value zero. Evaluating its fresh
step parameter $z$ at zero makes the successor $Sz$ have value one.
For case DI, the premise therefore supplies exactly the implication in
the displayed $D$ interpretation. For progress DI it supplies the same
implication, because its current hypothesis is evaluated at $z=0$.
The argument even applies to compound progress; the literal-equation
restriction is not needed for this invariant. The retained context is
unaffected by the fresh parameter.

In DE the fresh parameter $a$ also has value zero. Consequently the added
hypothesis $Q_yB(a,y)$ has exactly the truth value of its major $D_{x,y}B$.
Validity of the continuation proves validity of the conclusion. A
variable-target QE concludes the same truth value as its base premise;
a target-one QE is ordinary implication elimination. All primitive rules
preserve the invariant, so external induction on finite derivations
establishes soundness.

In (N1) the zero instance is $C(0,1,0)=1$, which is true, while the one
instance is $C(1,1,0)=1$, which says $0=1$ and is false. Hence (N1) is not
derivable under (R). In the unrestricted calculus, case DI with witness
$1$ proves it immediately from the computation equation
$C(Sz,1,1)=1$, ignoring the base hypothesis.
\end{proof}

\paragraph{General finite-range version.}
Fix $N\geq1$, let $W=\{0,\ldots,N\}$, allow QE only at variables and
numerals in $W$, and continue to restrict DI witnesses to zero or
variables. Evaluate free variables in $W$, keep ordinary numerical term
semantics (whose outputs may lie outside $W$), and interpret
\[
Q_yB := B(0)\Rightarrow\bigwedge_{j\in W} B(j),\qquad
D_{x,y}B := \bigvee_{i\in W}\left(B(i,0)\Rightarrow
                              \bigwedge_{j\in W}B(i,j)\right).
\]
A sequent is valid at every assignment of its free variables in $W$.
Case DI is valid because each positive $j\leq N$ is $S(j-1)$ with
$j-1\in W$. Progress DI is valid by the finite external induction from
zero to $N$. DE uses a witness in $W$ and its continuation's validity at
that assignment; QE uses a target in $W$. Equality transport is valid
because the recursively defined formula semantics respects numerical
equality even for substituted numerical values outside $W$.
The same primitive computations hold. Thus all rules preserve validity.
The unrestricted theorem
\[
D_{x,y}\bigl(C(y,N+1,x)=N+1\bigr)
\]
is false in this interpretation, since $1\in W$ but no possible witness
$i\in W$ equals $N+1$. Adding finitely many allowed numeral targets
therefore does not repair the simultaneous restriction.

\paragraph{Scope.}
This is an exact obstruction to identity-language theorem equivalence for
these simultaneous restrictions. It is not a proof that every differently
interpreted theory with restricted rules fails to be bi-interpretable with
iPRA. The special proof invariant for (N1) is stronger than the $N=1$
finite-range model and yields a smaller separating formula.

\paragraph{Compatibility with known-zero-instance Supply.}
The counterinterpretations remain sound for the rule
\[
\frac{\Gamma,H(w)\vdash A\qquad\Gamma\vdash H(0)}
     {\Gamma\vdash A},\qquad w\notin\operatorname{FV}(\Gamma,A).
\]
In the fixed-zero invariant, $H(w)$ and $H(0)$ have the same truth value.
For the finite-range invariant, extend any outer satisfying assignment
by $w=0$, which is permitted because $0\in W$. The second premise makes
$H(w)$ true at that extension; the first premise supplies $A$. Freshness
makes the end sequent unaffected by that extension. This also covers the
new lazy computation generator, whose separate soundness theorem proves
its generated equation's zero instance. No arbitrary-term naming rule
with a witness outside $W$ is included.

% END include

\clearpage
\section*{The inherited bridge, reproduced with attribution}
\addcontentsline{toc}{section}{The inherited bridge, reproduced with attribution}
The next five sections are copied unchanged from the generated-assumptions
archive at the pinned snapshot. That archive attributes the logical
Choice-DI basis, scalar predecessor construction, actual-equality
reflection, and primitive-recursion compiler to the earlier drafts.
Their exact source paths and SHA-256 identities are in
\file{evidence/source_manifest.json}. They are included so that the
bi-interpretation does not depend on an unavailable external proof.
None is claimed as a new result of this report. Their occurrences of
ordinary logic describe the interpretation theory; they do not add
primitives to the small native checker above.
% BEGIN appendices/inherited_choice_basis.tex
% New mathematical result: complementary introduction choices.
% Requires amsmath/amssymb/amsthm and theorem, lemma, proposition environments.
% The inherited arithmetic-to-iPRA bridge must be supplied separately.

\section{Complementary introductions with one local assumption}
\label{sec:choice-introduction}

The old paired introduction retains both the zero instance and an
induction instance when its matrix is an equation. Both together are
unnecessary. A certificate may choose either one, provided the induction
choice is restricted to a literal equation. This changes neither the
formula language nor the meaning of a formula. The omitted combination
is recovered by a short nested introduction--elimination construction.

\subsection{Complete native syntax and rule inventory}

The scalar numerical language is
\[
 t::=x\mid0\mid St\mid I(n,a;u.h).
\]
Only $h$ binds the state variable $u$; $n$ and $a$ do not.
The update may contain fixed numerical parameters. The computation
axioms in every retained context are
\[
 I(0,a;u.h)=a,\qquad
 I(Sn,a;u.h)=h[I(n,a;u.h)/u].
 \tag{I0, IS}\label{eq:choice-computations}
\]
The argument also applies unchanged to the full fixed-description PR
language with its projection, composition, and recursion equations.
No numerical variable ranges over function descriptions, and no uniform
evaluator is postulated.

Native formulas and reusable context clauses are
\[
 A::=(s=t)\mid\mathsf D_{x,y}A,\qquad
 \langle A;y.B(y)\rangle,\quad y\notin\operatorname{FV}(A).
\]
Bound variables are renamed before capture-free substitution.
The two variables $x,y$ in a paired formula are distinct.
An ordinary assumption $F$ is the constant clause
$\langle H;v.F\rangle$, where $v$ is absent from $F$ and $H$ is the
fixed computation theorem $I(0,0;u.0)=0$. Thus ordinary hypothesis
use is a clause Use preceded by this computation proof.

There are computation axioms and exactly the following inference
patterns:
\[
\frac{\Gamma\vdash A}{\Gamma\vdash B(t)}
\quad\mathsf{Use}
\quad(\langle A;y.B(y)\rangle\in\Gamma),
\]
\[
\frac{\Gamma\vdash\mathsf D_{x,y}B(x,y)\qquad
      \Gamma,\langle B(a,0);y.B(a,y)\rangle\vdash F}
     {\Gamma\vdash F}
\quad\mathsf{DE},
\qquad
\frac{\Gamma\vdash s=t\qquad\Gamma\vdash F(s)}
     {\Gamma\vdash F(t)}
\quad\mathsf{Tr},
\]
\[
\frac{\Gamma,B(t,\eta_\varepsilon(z))\vdash B(t,Sz)}
     {\Gamma\vdash\mathsf D_{x,y}B(x,y)}
\quad\mathsf{DI}_{\varepsilon},
\qquad
\eta_0(z)=0,\quad\eta_1(z)=z.
\tag{Choice DI}\label{eq:choice-di}
\]
Here $\varepsilon$ is a certificate choice in $\{0,1\}$, and
$\varepsilon=1$ is permitted only when $B$ has a literal equation
root. An equation beneath a binder is not a literal equation.
Call $\varepsilon=0$ the \emph{case choice} and $\varepsilon=1$
the \emph{progress choice}. Both retain exactly one ordinary
formula assumption. Compound matrices have only the case choice.

In DI the eigenparameter $z$ is fresh for $\Gamma$, the witness
$t$, and every fixed parameter of $B$. In DE the eigenparameter
$a$ is fresh for $\Gamma$, the conclusion $F$, and every fixed
parameter of $B$. In transport, $F(v)$ is any fixed native
formula template. Displayed local assumptions and clauses are
discharged. There are no primitive reflexivity, coverage,
separation, ordinary logical, or additional induction rules.

We call this calculus $\mathcal C$. The original comparison
calculus $\mathcal N$ has identical syntax, computation,
Use, DE, and transport, but its introduction is
\[
\frac{\Gamma,B(t,0),[B(t,z)]_{\mathrm{at}}\vdash B(t,Sz)}
     {\Gamma\vdash\mathsf D_{x,y}B(x,y)}.
\tag{Old DI}\label{eq:choice-old-di}
\]
The bracket supplies $B(t,z)$ exactly at a literal equation.

\subsection{Exact equivalence and the unused-witness construction}

\begin{lemma}[Structural admissibility]
\label{lem:choice-structural}
Weakening, exchange, contraction, renaming, ordinary-assumption cut,
and capture-free numerical substitution are admissible in $\mathcal C$.
\end{lemma}
\begin{proof}
Induct on derivations. Rename local eigenparameters before inserting
proofs or terms. Numerical substitution preserves a formula's root
constructor, so an allowed progress choice remains allowed.
The selected choice does not change under substitution or cut.
For ordinary-assumption cut, replace each use of its constant clause
by the supplied proof. Computation schemas are closed under
capture-free numerical substitution. These facts verify every case.
\end{proof}

\begin{theorem}[Identity-language equivalence]
\label{thm:choice-equivalence}
For every retained clause context $\Gamma$ and native formula $F$,
\[
 \Gamma\vdash_{\mathcal C}F
 \quad\Longleftrightarrow\quad
 \Gamma\vdash_{\mathcal N}F.
\]
The transformations are effective on finite derivations and leave
all numerical and formula syntax unchanged.
\end{theorem}
\begin{proof}
Each new introduction is an instance of the old one after weakening
by the omitted local assumption. For the case choice at an equation,
add the induction instance. For the progress choice, add the zero
instance. Compound case introduction is already the old rule.
Thus every new derivation translates to an old one.

For the reverse direction, transform an old derivation recursively.
Only its atomic DI case needs a construction. Suppose its matrix is
the equation $E(x,y)$ and its transformed premise is
\[
 \pi:\quad\Gamma,E(t,0),E(t,z)\vdash E(t,Sz).
\]
The outer new inference will be case DI with witness $t$. It is
therefore enough to prove $E(t,Sz)$ under
$\Gamma,E(t,0)$.

Rename the eigenvariable in $\pi$ to a fresh $v$. Choose a new
variable $u$ absent from $E(t,y)$ and all retained data. Under
$\Gamma,E(t,0)$, progress DI with witness zero introduces
\[
             \mathsf D_{u,y}E(t,y)
\]
from exactly the renamed premise $\pi[v/z]$.
The bound coordinate $u$ is deliberately unused.

Apply DE to this paired formula. Because $u$ is unused, the opened
clause is exactly
\[
               \langle E(t,0);y.E(t,y)\rangle;
\]
it contains no new existential witness. The retained ordinary
assumption $E(t,0)$ supplies its antecedent. Use at $Sz$ gives
$E(t,Sz)$, and DE closes. Finally the outer case DI introduces
the original $\mathsf D_{x,y}E(x,y)$ with witness $t$.
The original premise proof $\pi$ has been inserted once.

An old compound introduction already is the new case choice.
All other rule cases are unchanged. This completes the induction
on the original derivation.
\end{proof}

The unused inner coordinate is essential. Introducing
$\mathsf D_{x,y}E(x,y)$ internally and then eliminating it would
give an arbitrary fresh witness, not necessarily the original
$t$. The construction instead keeps $t$ as a fixed parameter
of $E(t,y)$ and binds only the unused coordinate $u$.

\begin{proposition}[Linear bound on inference nodes]
\label{prop:choice-node-bound}
Count each computation axiom, Use, DI, DE, and transport as one
inference node, expanding ordinary-assumption use as a computation
of $H$ followed by Use. If an old derivation has $N$ nodes and
$r$ atomic DI nodes, the displayed transformation has exactly
\[
                   N+5r\leq 6N
\]
nodes.
\end{proposition}
\begin{proof}
Each atomic old DI node is replaced by six nodes: outer case DI,
inner progress DI, DE, Use of the opened clause, Use of the
retained base assumption, and the computation of $H$ needed for
that last ordinary-assumption use. Its old premise subtree
occurs once. All other inference nodes are copied once.
\end{proof}

This counts inference-tree nodes. It is not a bound on serialized
JSON bytes, formula sizes, checker runtime, or the number of
consumed packet cells when hypothesis closure is implicit.
In the formula-context Q/D presentation, the corresponding
replacement has outer case DI, inner progress DI, DE, QE,
and two Hyp nodes, giving the same six-node count.

\subsection{Induction and numerical cases separate cleanly}

\begin{lemma}[Contextual equation induction]
\label{lem:choice-ei}
The progress choice, DE, and Use derive
\[
\frac{\Gamma\vdash E(0)\qquad
      \Gamma,E(z)\vdash E(Sz)}
     {\Gamma\vdash E(t)}
\tag{EI}\label{eq:choice-ei}
\]
for every equation $E$, with the usual freshness condition.
\end{lemma}
\begin{proof}
Choose an unused witness coordinate $u$. Progress DI introduces
$\mathsf D_{u,y}E(y)$ from the given step. DE opens the clause
$\langle E(0);y.E(y)\rangle$. Supply its antecedent using the given
base proof and use the clause at $t$. Close DE.
\end{proof}

\begin{lemma}[Cases, equality, and binder congruence]
\label{lem:choice-preinduction}
The case choice derives, for arbitrary native formulas,
\[
\frac{\Gamma\vdash F(0)\qquad\Gamma\vdash F(Sz)}
     {\Gamma\vdash F(t)}
\tag{NC}\label{eq:choice-nc}
\]
with the usual freshness condition. Equality reflexivity,
symmetry, transitivity, numerical congruence, and uniform
matrix interderivability under $\mathsf D$ are derivable
without using the progress choice.
\end{lemma}
\begin{proof}
For cases, use case DI with an unused witness coordinate and
matrix $F(y)$, ignoring its retained base in the successor
premise. DE opens $\langle F(0);y.F(y)\rangle$; the given zero
proof and Use at $t$ finish.

For reflexivity put $r=I(0,t;u.0)$. Use the computation equation
$r=t$ as both transport premises with motive $v=t$ to obtain
$t=t$. The usual explicit equality motives give symmetry,
transitivity, and numerical congruence.

For binder congruence, open the old paired formula at a fresh
witness $a$ and introduce the new one by case DI at $a$.
Convert its retained base to the old base, use the opened
old clause at the successor, and convert back. Ordinary cut
closes the construction, and the reverse conversion gives
the converse. No induction is used.
\end{proof}

The primitive progress choice is still restricted to equations.
Allowing it at arbitrary compound matrices would, by the
argument for \eqref{eq:choice-ei}, supply induction for
arbitrary formulas and change the intended theory.

\subsection{The entire ordinary logical interface}
\label{subsec:choice-logic}

Here is a predecessor-free reconstruction of the ordinary logical
operations. It makes explicit that the case choice supplies the
logical interface independently of progress induction.

Let $C(v,a,b):=I(v,a;u.b)$ with $u$ unused. Computation gives
$C(0,a,b)=a$ and $C(Sv,a,b)=b$. Define the gate on fixed finite
formula syntax by
\[
 G_v(s=t):=(C(v,0,s)=C(v,0,t)),\qquad
 G_v(\mathsf D_{x,y}B):=\mathsf D_{x,y}G_vB,
\]
renaming bound variables away from the control. Equality,
binder congruence, and external structural induction give
\[
 \vdash G_0A,\qquad G_{Sv}A\dashv\vdash A,\qquad
 A\dashv\vdash\mathsf D_{u,y}G_yA
 \quad(u,y\notin\operatorname{FV}(A)).
\tag{Gate and padding}\label{eq:choice-gate}
\]
For padding, case DI proves the forward direction. DE, the
zero-gate theorem, and Use at $1$ prove the reverse direction.

Define a selector between two arbitrary formulas by
\[
 M_v(s=t,u=w):=(C(v,s,u)=C(v,t,w))
\]
and, with aligned fresh paired variables, by
\[
\begin{aligned}
 M_v(\mathsf D_{x,y}A,\mathsf D_{x,y}B)
   &:=\mathsf D_{x,y}M_v(A,B),\\
 M_v(E,\mathsf D_{x,y}B)
   &:=\mathsf D_{x,y}M_v(G_yE,B),\\
 M_v(\mathsf D_{x,y}A,E)
   &:=\mathsf D_{x,y}M_v(A,G_yE),
\end{aligned}
\]
where $E$ is an equation. Every call decreases maximum input
binder depth. Padding and congruence prove
$M_0(A,B)\dashv\vdash A$ and
$M_{Sv}(A,B)\dashv\vdash B$.

Put $\delta(v)=C(v,1,0)$ and $p(z)=I(z,0;u.\delta(u))$.
Only the fixed computations $p(1)=1$ and $p(2)=0$ are needed.
The full logical macro inventory is
\[
\begin{aligned}
 A\Rightarrow_*B&:=\mathsf D_{u,z}M_z(A,B),\\
 A\wedge_*B&:=\mathsf D_{u,z}G_zM_{p(z)}(A,B),\\
 \exists_*x\,A(x)&:=\mathsf D_{x,y}G_yA(x),\\
 A\vee_*B&:=\exists_*x\,M_x(A,B),\\
 \forall_*x\,A(x)&:=
        A(0)\wedge_*\mathsf D_{u,z}G_zA(z),\\
 \bot_*&:=(0=1).
\end{aligned}
\tag{Logical macros}\label{eq:choice-macros}
\]
Newly bound variables are fresh; $u$ is unused.

Implication introduction obtains its assumption from the zero
selector and supplies the consequent at successors, then uses
case DI. Elimination opens its clause and uses it at $1$.
Existential introduction uses its witness and the successor
gate; elimination opens at a fresh witness, supplies the
zero-gate theorem, and uses the clause at $1$.

Disjunction introduction uses witnesses $0$ and $1$.
For elimination, discharge the two branch proofs as
$A\Rightarrow_*F$ and $B\Rightarrow_*F$. NC applied to
$M_z(A,B)\Rightarrow_*F$ gives this implication at every
control: the two endpoint cases follow from the branch
implications. Existential elimination provides $M_a(A,B)$
and the implication gives $F$.

Given $A$ and $B$, selector endpoints and NC give $M_v(A,B)$
at an arbitrary control. Substitute $v=p(Sz)$ and use the
successor gate for conjunction introduction. For projections,
open its zero-gated clause and specialize at $1$ and $2$;
the two fixed computations give $B$ and $A$, respectively.

These rules prove existential numerical coverage:
\[
     (x=0)\vee_*\exists_*v\,(x=Sv).
\tag{Coverage}\label{eq:choice-coverage}
\]
Use NC. Zero takes the left injection of reflexivity, and a
successor takes the right injection with the evident witness.
Coverage gives equality cases even when the retained context
contains $x$; its elimination adds $x=0$ or $x=Sv$ rather
than silently changing that retained context.

Universal introduction follows from a fresh-parameter proof
by substituting zero and successors and forming the two
conjuncts. For elimination, take its $A(0)$ conjunct and
open the paired conjunct. Its base is a zero-gate theorem.
Coverage of the target chooses the zero instance from
$A(0)$ or a successor instance from the opened clause
and gate removal.

Finally, under $0=1$, computation and transport give
$s=C(0,s,t)=C(1,s,t)=t$ for arbitrary numerical terms.
Every equation follows, and external induction on formula
syntax with case DI gives every native formula. This
supplies falsity elimination. No excluded middle for
quantified formulas has entered the construction.

\subsection{Whole-domain realizations and both formula composites}

For a standard first-order reading, let $T_{\rm eq}$ have
intuitionistic equality logic, the same numerical signature
and computation equations, numerical coverage, separation
$(0=1)\to\bot$, and contextual equation induction.
Its interpretation of the native syntax is
\[
 (s=t)^\circ=(s=t),\qquad
 (\mathsf D_{x,y}B)^\circ
   =\exists x\bigl(B^\circ(x,0)\to\forall y B^\circ(x,y)\bigr),
\]
\[
 \langle A;y.B\rangle^\circ
       =A^\circ\to\forall y B^\circ.
\]
The domain is the entire numerical domain and interpreted
equality is actual equality.

Case DI is sound in this realization by numerical coverage.
Progress DI is sound by contextual equation induction under
the conditional base. Use, DE, transport, and computation
have their displayed ordinary meanings. Conversely, the
native macros derive every ordinary logical rule,
\eqref{eq:choice-ei} derives equation induction, and
\eqref{eq:choice-coverage} derives coverage. Separation
translates to the identity implication
$(0=1)\Rightarrow_*(0=1)$.

For the essential paired formula composite,
\[
 \mathsf D_{x,y}B
 \dashv\vdash
 \exists_*x\bigl(B(x,0)\Rightarrow_*\forall_*y\,B(x,y)\bigr).
\tag{Paired composite}\label{eq:choice-paired-composite}
\]
Forwards, DE opens its clause at a fresh $a$; under $B(a,0)$
that clause supplies $B(a,y)$ at a fresh arbitrary $y$.
Introduce the universal, implication, and existential
macros and close DE. Backwards, eliminate the existential
at a fresh $a$. Introduce the paired formula with witness
$a$ by case DI: the retained base, supplied implication,
and universal elimination at $Sz$ give its successor premise.

The ordinary interpretations of the six logical macros
have their ordinary round trips. For conjunction, its true
base and the instances $1,2$ recover both operands;
the converse uses selector cases. The universal has all
successor instances in its paired conjunct and zero in
its first conjunct; coverage gives exactly the ordinary
universal. Implication, disjunction, and existential
quantification use their selector or gate endpoints.
Falsity uses the explicit ordinary separation bridge.
External induction on finite formulas, together with
\eqref{eq:choice-paired-composite}, gives both formula
composites. Clauses correspond in both directions to
$A\Rightarrow_*\forall_*yB$ by the derived logical rules.

The numerical composites are literal identities with
internally provable graph $x=y$. No standard-numeral
restriction, retract, quotient, or equality reflection
assumption is being inserted at this stage. The separately
supplied arithmetic bridge identifies $T_{\rm eq}$ with
full iPRA; the present change preserves that bridge
unchanged because Theorem~\ref{thm:choice-equivalence}
preserves every native sequent.

\subsection{What the theorem does and does not claim}

The new inference interface retains one ordinary assumption
at every paired introduction. It preserves the entire
consequence relation, not just equations or closed
theorems. Consequently the established full-iPRA
interpretations and equiconsistency transfer unchanged.

This result is different from retaining only progress
introduction and strengthening closure elimination.
Here the original elimination rule is retained, and
the case and progress introductions coexist. It is also
different from allowing arbitrary compound progress,
which would supply stronger induction.

A complete-checker comparison must still count the
certificate choice, its equation-root restriction,
substitution, context handling, computation, and the
public input/output interface. The mathematical theorem
alone is not a byte-count or minimality theorem.
The inference-node bound is distinct from source-size
measurements and from serialized certificate size.
The proof is a written universal proof; finite regression
tests do not constitute a Lean/Rocq formalization of it.

% END include

% BEGIN appendices/scalar_predecessor.tex
% Adapted with attribution from the pinned unparenthesized one-hole archive.
\section{Internal scalar predecessor bootstrap}
\label{sec:hybrid-bootstrap}

This scalar bootstrap is inherited from the supplied scalar development~\cite{erasure} and is reproduced to keep the scalar checker self-contained. It is not a new contribution of this paper. The proof uses a factor map fixing odd states and a saturated-iteration construction of predecessor.

The available arithmetic terms are $0$, $S$, and
\[
 I(n,a;u.h),\qquad
 I(0,a;u.h)=a,\qquad
 I(Sn,a;u.h)=h[I(n,a;u.h)/u].
\]
The update binds only $u$.  All other variables in it are fixed numerical
parameters.  We use equality reasoning, the already reconstructed
zero/successor cases, and induction on a single literal equation.
Every induction below keeps its other parameters fixed.  Predecessor,
successor injectivity, pairing, and induction on arbitrary
quantifier-free formulas are not assumptions of the argument.

\subsection{Elementary terms and two finite cycles}

Write $1=S0$ and use larger numerals as successor abbreviations.  Define
\[
 C(n,a,b)=I(n,a;u.b),\qquad
 Z(n)=C(n,1,0),\qquad
 a+b=I(b,a;u.Su).
\]
The update variable of $C$ is absent from its body.  Consequently
\[
 C(0,a,b)=a,\qquad C(Sn,a,b)=b.
\]
Numerical cases show that $Z(n)$ is $0$ or $1$.
The usual commutative-monoid equations for addition follow by literal
equation induction.  For example, $0+b=b$ and $Sa+b=S(a+b)$ are
proved by induction on $b$; these yield commutativity by a further
equation induction, and associativity follows by induction on the last
summand.  We use these equations to reorder finite sums.  The notation
$2a$ means $a+a$ and $3a$ means $a+a+a$; no multiplication operation
is needed for this bootstrap.

Put
\[
 N(n)=I(n,0;u.Z(u)),\qquad
 j(u)=Z(u)+2N(u),\qquad
 M(n)=I(n,0;u.j(u)).
\]
The range of $N$ is contained in $\{0,1\}$: its zero value is $0$,
and its successor value is a value of $Z$.  This uses numerical cases
on its counter, not induction on a disjunction.  Hence
\begin{equation}
 N(n)+Z(N(n))=1,\qquad
 Z(Z(N(n)))=N(n),\qquad N(n+2)=N(n).
 \label{eq:hybrid-parity}
\end{equation}
Here the last equation follows from two successor computations and
the preceding equation.  Literal-equation inductions give
\begin{equation}
 N(2n)=0,\qquad N(2n+1)=1,\qquad N(3n)=N(n).
 \label{eq:hybrid-parity-scaling}
\end{equation}
For the first equation the successor step adds two to the argument of
$N$; the second follows by one computation.  In the third, increasing
$n$ adds three to the left argument and therefore applies $Z$ once to
both sides of the induction equation.

Finite computations now give
\[
 j(0)=1,\qquad j(1)=2,\qquad j(2)=0.
\]
The proof below needs just these three values and
$M(Sn)=j(M(n))$; it does not assume a division algorithm or a theorem
characterizing $M$ as remainder modulo three.

\subsection{A ramp with a global scaling identity}

Define
\[
 G(n)=I(n,0;u.u+M(u)+1).
\]
The central invariant is
\begin{equation}
 M(G(n))=N(n).
 \label{eq:hybrid-ramp-invariant}
\end{equation}
Its zero case is immediate.  For its successor step, retain the
single induction equation and split the Boolean value $N(n)$.
If it is $0$, the next $G$ state is $G(n)+1$; its $M$ value is
$j(0)=1$, equal to $N(Sn)$.  If it is $1$, the next $G$ state
is $G(n)+2$; two computations of $M$ change its value from $1$
to $2$ and then to $0$, again equal to $N(Sn)$.  This proves
the invariant by literal-equation induction.

It follows that
\[
 G(Sn)=G(n)+N(n)+1,\qquad G(n+2)=G(n)+3.
\]
Equation induction therefore yields
\begin{equation}
 G(2n)=3n.
 \label{eq:hybrid-ramp-double}
\end{equation}
A separate literal-equation induction proves
\begin{equation}
 G(3n)=3G(n)+N(n).
 \label{eq:hybrid-ramp-triple}
\end{equation}
For completeness, its successor calculation is
\begin{align*}
 G(3Sn)
 &=G(3n)+N(3n)+4
   =3G(n)+2N(n)+4,\\
 3G(Sn)+N(Sn)
 &=3G(n)+3N(n)+3+Z(N(n)).
\end{align*}
The expressions agree in both Boolean cases for $N(n)$.
The base is $0=0$.  In particular,
\eqref{eq:hybrid-ramp-triple} was proved internally; it was not
inferred from an intended floor-function interpretation of $G$.

Set
\[
 H(u)=C(N(u),G(u),u).
\]
The crucial choice is to retain $u$ when its parity is $1$.
The preceding equations give
\begin{equation}
 H(2u)=3u,\qquad H(1)=1,\qquad H(3u)=3H(u).
 \label{eq:hybrid-factor-map}
\end{equation}
To check the last identity, use $N(3u)=N(u)$ and split the Boolean
$N(u)$.  At $0$, both sides are $3G(u)$ by
\eqref{eq:hybrid-ramp-triple}; at $1$, both are $3u$.
This argument does not require that an arbitrary number has already
been represented as $2v$ or $2v+1$.

\subsection{A comparator from scalar iteration}

For every fixed update $h$, equation induction on $n$, keeping $m,a$
and all parameters fixed, proves the semigroup identity
\[
 I(m+n,a;h)=I(n,I(m,a;h);h).
\]
The base is the zero computation and the step applies $h$ to the
induction equation. This identity will also be used in the later
pairing construction.


For a fixed update $h$, equation induction on $y$, with its initializer
$a$ and all parameters of $h$ fixed, proves the rotation identity
\begin{equation}
 I(Sy,a;h)=I(y,h(a);h).
 \label{eq:hybrid-rotation}
\end{equation}
Both sides have base $h(a)$, and their successor equations apply the
same $h$ to the previous equality.  Here and below $I(y,a;h)$ is
an abbreviation for one fixed finite scalar term, not quantification
over functions.

By \eqref{eq:hybrid-factor-map}, another equation induction gives
\begin{equation}
 I(y,3a;H)=3I(y,a;H).
 \label{eq:hybrid-iterator-scaling}
\end{equation}
Its base is $3a=3a$.  In the step the fixed-parameter induction
equation rewrites the new left side as
\[
 H(3I(y,a;H))=3H(I(y,a;H)).
\]
The initializer parameter stays $a$ throughout this induction.
Equation induction also gives $I(y,1;H)=1$, since $H(1)=1$.

Define
\[
 E(x)=I(x,1;u.2u),\qquad
 L(x,y)=N\bigl(I(y,E(x);H)\bigr).
\]
Then the following entire recursion specification is provable:
\begin{equation}
 L(0,y)=1,\qquad L(Sx,0)=0,\qquad L(Sx,Sy)=L(x,y).
 \label{eq:hybrid-comparator}
\end{equation}
The first equation uses $E(0)=1$ and the fixed-point identity for
$H$.  The second uses $E(Sx)=2E(x)$ and $N(2a)=0$.
For the third,
\begin{align*}
 L(Sx,Sy)
 &=N\bigl(I(Sy,2E(x);H)\bigr)\\
 &=N\bigl(I(y,H(2E(x));H)\bigr)\\
 &=N\bigl(I(y,3E(x);H)\bigr)\\
 &=N\bigl(3I(y,E(x);H)\bigr)\\
 &=N\bigl(I(y,E(x);H)\bigr)=L(x,y).
\end{align*}
The successive steps use rotation, $H(2a)=3a$, iterator scaling,
and $N(3a)=N(a)$.  Thus $L$ has the comparison equations at all
arguments internally.  Its standard-model meaning as the indicator
of $x\le y$ is motivation, not a premise of the following argument.

\subsection{Saturated iteration gives predecessor}

We now need only the three equations
\eqref{eq:hybrid-comparator}.  Define
\[
 B_m(u)=C(L(Su,m),u,Su),\qquad
 W(t,m)=I(t,0;u.B_m(u)).
\]
The subscript in $B_m$ records a fixed parameter; $B_m$ is an open
term abbreviation, not a function variable.

Numerical cases on $e$ prove the selector equation
\[
 C(e,Sa,Sb)=S C(e,a,b).
\]
Equivalently one may use equation induction on $e$, whose induction
hypothesis is unused.  Therefore
\begin{align}
 B_{Sm}(Su)
 &=C(L(SSu,Sm),Su,SSu)\notag\\
 &=C(L(Su,m),Su,SSu)=S B_m(u),
 \label{eq:hybrid-conjugacy}\\
 B_{Sm}(0)&=C(L(S0,Sm),0,S0)=C(L(0,m),0,S0)=S0.
 \label{eq:hybrid-conjugacy-base}
\end{align}

Literal-equation induction on $t$, with $m$ fixed, proves
\begin{equation}
 W(St,Sm)=S W(t,m).
 \label{eq:hybrid-saturated-shift}
\end{equation}
The base is \eqref{eq:hybrid-conjugacy-base}.  The successor step is
\begin{align*}
 W(SSt,Sm)
 &=B_{Sm}(W(St,Sm))\\
 &=B_{Sm}(S W(t,m))\\
 &=S B_m(W(t,m))=S W(St,m).
\end{align*}
Only the second line uses the current induction hypothesis, at the
same fixed $m$.

Two subsequent equation inductions give
\begin{equation}
 L(Sn,n)=0,\qquad W(n,n)=n.
 \label{eq:hybrid-diagonals}
\end{equation}
The first uses the comparator recurrence.  The second uses
\eqref{eq:hybrid-saturated-shift} in its successor step.  The
substitution $t=m=n$ is made into that already proved universally
parameterized equation; no parameter of an active induction hypothesis
is changed.  Consequently
\begin{equation}
 W(Sn,n)=B_n(n)=C(L(Sn,n),n,Sn)=n.
 \label{eq:hybrid-extra-step}
\end{equation}

Finally define the finite scalar term
\begin{equation}
 P(m)=I\bigl(m,0;u.C(L(SSu,m),u,Su)\bigr).
 \label{eq:hybrid-predecessor-term}
\end{equation}
Its zero equation is immediate.  To justify its successor equation
without assuming a separate extensionality rule for iterator bodies,
put
\[
 V(t,m)=I\bigl(t,0;u.C(L(SSu,Sm),u,Su)\bigr).
\]
The comparator recurrence identifies this update pointwise with
$B_m(u)$.  Equation induction on $t$, with $m$ fixed, gives
$V(t,m)=W(t,m)$: their base values are $0$, and the successor step
uses the pointwise update equation and then the induction equation.
Now \eqref{eq:hybrid-extra-step} yields
\[
 \boxed{\qquad P(0)=0,\qquad P(Sm)=V(Sm,m)=W(Sm,m)=m.\qquad}
\]

\subsection{Dependency boundary}

All terms in this section are built from $0,S,I$ by finite substitution.
All inductions have a single equation as their invariant.  Numerical
cases and equality transport are the already reconstructed logical
resources.  Successor injectivity may now be obtained by applying $P$
to an assumed equation $Sx=Sy$.  Equality tests, reflection of
quantifier-free formulas, pairing, and the recovery of general
primitive recursion can then be developed in the established order.
They were not used in order to manufacture predecessor.

This argument is a written internal derivation in the stated theory.
The finite implementation checks do not constitute a proof-assistant
formalization of this entire metatheorem.

% END include

% BEGIN appendices/equality_reflection.tex
% Standalone section.  Requires only amsmath and amssymb in the parent document.
% No theorem environments, macros, or bibliography keys are assumed.
% All arithmetic claims in this section are proved here.

\section{From literal-equation induction to full intuitionistic PRA}
\label{sec:eic-completion}\label{app:arithmetic}

This section supplies an unconditional completion theorem. It is adapted
with attribution from the current forward-left draft~\cite{forward}. In the
scalar application used here, every function occurring in the proof is
already a scalar term: predecessor is supplied by
Appendix~\ref{sec:hybrid-bootstrap}, the selector is
$C(n,a,b)=I(n,a;u.b)$, addition is $x+y=I(y,x;u.Su)$, and
truncated subtraction is $d(x,y)=I(y,x;u.Pu)$. All other functions below
are finite compositions of these. Thus this proof does not presuppose
that arbitrary primitive recursion has already been compiled. The QFI
conclusion first applies to the scalar language; the general PR compiler
comes afterward.  Literal-equation
induction, together with the primitive-recursive defining equations,
predecessor coverage, and separation of $0$ and $1$, proves induction for every
quantifier-free formula.  In particular, the proof does not assume a faithful
numerical equality test, stability or decidability of equality, successor
uniqueness, or any commutation law for an iterator.  We construct the equality
test and prove its reflection property for arbitrary elements of the
numerical domain.

\subsection{The exact hypotheses}

Work in ordinary intuitionistic first-order logic with equality and genuine
falsity $\bot$.  The language has one numerical sort and a symbol for each
fixed, finite primitive-recursive description.  Distinct descriptions are
not identified merely because they compute the same standard function.
Let $\mathsf{PRE}$ contain the defining equations for zero, successor,
projections, composition, and primitive recursion.  In particular, choose
descriptions $P$ and $C$ with equations
\begin{align}
 P(0)&=0, & P(Sx)&=x, \label{eq:eic-predecessor}\\
 C(0,a,b)&=a, & C(Sx,a,b)&=b. \label{eq:eic-selector}
\end{align}
Write $1=S0$.  In addition to $\mathsf{PRE}$, assume
\begin{align}
 \mathsf{Cov}:&\qquad
       \forall x\,\bigl(x=0\lor x=S(Px)\bigr),
       \label{eq:eic-coverage}\\
 \mathsf{Sep}:&\qquad 0=1\longrightarrow\bot,
       \label{eq:eic-separation}
\end{align}
and induction for each \emph{literal equation} $E(n,\vec p)$:
\begin{equation}
 \forall\vec p\,
 \left[
   \left(E(0,\vec p)\land
       \forall n\bigl(E(n,\vec p)\longrightarrow E(Sn,\vec p)\bigr)
   \right)
   \longrightarrow\forall n\,E(n,\vec p)
 \right].
 \label{eq:eic-induction}
\end{equation}
Here the schematic expression $E$ is a single equality of numerical
terms; the displayed universal quantifiers range only over numbers.
Every description and every instance of the schema is finite external
syntax.

The axiom formulation \eqref{eq:eic-induction} and the usual contextual
induction rule are interchangeable.  For example, if
\[
 \Gamma\vdash E(0,\vec p),
 \qquad
 \Gamma,E(z,\vec p)\vdash E(Sz,\vec p),
\]
where $z$ is absent from $\Gamma$ and the passive parameters, discharge the
second premise's equation and universally introduce $z$.  The axiom then
gives $\Gamma\vdash\forall n E(n,\vec p)$, which can be instantiated at any
target term.  Conversely, assume the antecedent of the axiom and apply the
contextual rule with fresh variables.  This uses the ordinary intuitionistic
logical rules and permits arbitrary, including quantified, formulas in
$\Gamma$.

\subsection{Constructor separation and two elementary tests}

Define the sign description and the countdown description by
\begin{equation}
 \operatorname{sg}(u)=C(u,0,1),
 \qquad
 K(u)=C(u,1,Pu).
 \label{eq:eic-sign-countdown}
\end{equation}
Their defining equations imply
\begin{align}
 \operatorname{sg}(0)&=0,
 &\operatorname{sg}(Sx)&=1,\label{eq:eic-sign-computation}\\
 K(0)&=1,
 &K(Sx)&=x.\label{eq:eic-countdown-computation}
\end{align}

\paragraph{Constructor laws.}
Successor is injective: applying $P$ to $Sx=Sy$ gives $x=y$ by
\eqref{eq:eic-predecessor}.  Successor never equals zero.  Indeed,
$Sx=0$ gives $1=0$ by applying $\operatorname{sg}$ and using
\eqref{eq:eic-sign-computation}; symmetry and
\eqref{eq:eic-separation} give $\bot$.  Consequently $1\ne0$, where
$A\ne B$ abbreviates $A=B\to\bot$.  None of these arguments uses
induction.

\paragraph{The sign test.}
For every $u$,
\begin{align}
 &\operatorname{sg}(u)=0\lor\operatorname{sg}(u)=1,
       \label{eq:eic-sign-range}\\
 &\operatorname{sg}(u)=0\quad\longleftrightarrow\quad u=0,
       \label{eq:eic-sign-zero}\\
 &\operatorname{sg}(u)=1\quad\longleftrightarrow\quad u\ne0.
       \label{eq:eic-sign-one}
\end{align}
To prove the range assertion, split \eqref{eq:eic-coverage}.  The branch
$u=0$ gives sign $0$, and the branch $u=S(Pu)$ gives sign $1$.
The implication $u=0\to\operatorname{sg}(u)=0$ is computation.
For its converse, the successor branch of coverage would make
$\operatorname{sg}(u)=0$ imply $1=0$, a contradiction; in that branch
explosion gives $u=0$ as well.  Similarly, if
$\operatorname{sg}(u)=1$ and $u=0$, computation gives $0=1$.
Conversely, $u\ne0$ excludes the zero branch of coverage, and the
successor branch computes the sign as $1$.

Coverage and \eqref{eq:eic-selector} also give, for all $u,v$,
\begin{equation}
 C(u,v,v)=v.
 \label{eq:eic-selector-idempotence}
\end{equation}
Both coverage branches compute to $v$, so this identity again needs no
induction.

\paragraph{The countdown test and the notation $W$.}
For every $u$,
\begin{equation}
 K(u)=0\quad\longleftrightarrow\quad u=1.
 \label{eq:eic-countdown-zero}
\end{equation}
If $u=1=S0$, \eqref{eq:eic-countdown-computation} gives $K(u)=0$.
Conversely, assume $K(u)=0$ and split coverage.  If $u=0$, then
$K(u)=1$, contradicting separation, and explosion gives $u=1$.
If $u=S(Pu)$, computation gives $K(u)=Pu$; hence $Pu=0$ and
$u=S0=1$.

If a presentation uses a description named $W$ with defining equations
\[
 W(0)=1,\qquad W(Sx)=x,
\]
then it proves $K(u)=W(u)$ for every $u$ already from coverage.
In the zero branch both sides are $1$; in the branch $u=S(Pu)$ both
sides are $Pu$.  Thus $K$ and $W$ are provably equal on the whole domain,
even if they are distinct descriptions.  This equality does not appeal to
standard-model agreement or to uniqueness of an arbitrary recurrence.

\subsection{Four elementary uses of equation induction}

Choose descriptions for addition and truncated subtraction, with recursion
on the second argument:
\begin{align}
 x+0&=x, &x+Sy&=S(x+y),\label{eq:eic-addition}\\
 d(x,0)&=x, &d(x,Sy)&=P(d(x,y)).\label{eq:eic-subtraction}
\end{align}
Using these equations and \eqref{eq:eic-induction}, we now derive
\begin{align}
 0+y&=y,\label{eq:eic-left-zero}\\
 (Sx)+y&=S(x+y),\label{eq:eic-left-successor}\\
 d(Sx,Sy)&=d(x,y),\label{eq:eic-cancel-successors}\\
 d(Sa,a)&=1.\label{eq:eic-diagonal}
\end{align}
Each induction below has a single equation as its induction formula.

For \eqref{eq:eic-left-zero}, induct on $y$.  The base is $0+0=0$.
The step is $0+Sy=S(0+y)=Sy$.

For \eqref{eq:eic-left-successor}, hold $x$ fixed and induct on $y$.
At zero both sides are $Sx$.  At a successor the induction hypothesis
gives
\[
 (Sx)+Sy=S\bigl((Sx)+y\bigr)
       =S\bigl(S(x+y)\bigr)=S(x+Sy).
\]
Combining \eqref{eq:eic-left-successor} with the defining equation for
addition gives the useful rearrangement
\begin{equation}
 (Sy)+r=y+Sr.
 \label{eq:eic-shift-successor}
\end{equation}
Indeed, both sides are $S(y+r)$.

For \eqref{eq:eic-cancel-successors}, hold $x$ fixed and induct on $y$.
The base computes as
\[
 d(Sx,S0)=P(d(Sx,0))=P(Sx)=x=d(x,0).
\]
Applying $P$ to the induction equation gives the step:
\[
 d(Sx,S(Sy))=P(d(Sx,Sy))=P(d(x,y))=d(x,Sy).
\]

Finally, induct on $a$ for \eqref{eq:eic-diagonal}.  The base is
$d(S0,0)=S0=1$.  At a successor,
\[
 d(S(Sa),Sa)=d(Sa,a)=1,
\]
using \eqref{eq:eic-cancel-successors} and the induction hypothesis.

\subsection{An equation invariant that reflects actual equality}

The remaining arithmetic ingredient is the following equation, valid for
all $x,y$:
\begin{equation}
 C\bigl(\operatorname{sg}(d(x,y)),\ x,\ y+d(x,y)\bigr)=x.
 \label{eq:eic-residual-invariant}
\end{equation}
Its meaning is that a nonzero subtraction residual reconstructs the
original argument by addition.  The selector expresses this assertion as
one literal equation, so equation induction applies directly.

Hold $x$ fixed and induct on $y$.  At $y=0$, subtraction gives residual
$x$, and \eqref{eq:eic-left-zero} gives $0+x=x$.  The desired equation
is consequently $C(\operatorname{sg}(x),x,x)=x$, which is
\eqref{eq:eic-selector-idempotence}.

For the step, assume \eqref{eq:eic-residual-invariant} at $y$ and
abbreviate
\[
 r=d(x,y),\qquad r'=Pr=d(x,Sy).
\]
These are local term abbreviations, not additional function symbols.
Split coverage for $r'$.

If $r'=0$, then $\operatorname{sg}(r')=0$, and the successor instance
of \eqref{eq:eic-residual-invariant} reduces immediately to $x=x$.

In the other branch, $r'=S(Pr')$, so constructor separation gives
$r'\ne0$.  Split coverage for $r$.  Its zero branch would imply
$r'=P0=0$, a contradiction, from which the desired step equation follows.
In its successor branch,
\[
 r=S(Pr)=Sr',\qquad \operatorname{sg}(r)=1.
\]
The induction equation therefore computes to $y+r=x$.  Meanwhile,
\eqref{eq:eic-shift-successor} gives
\[
 (Sy)+r'=y+Sr'=y+r=x.
\]
Since $r'$ is a successor, $\operatorname{sg}(r')=1$, so this is
exactly the branch selected in the successor instance of
\eqref{eq:eic-residual-invariant}.  This completes the equation induction.
All case distinctions came from coverage; no induction on an implication
or a disjunction was used.

\paragraph{Reflection at residual $1$.}
We have, for arbitrary $a,b$,
\begin{equation}
 d(Sa,b)=1\quad\longleftrightarrow\quad a=b.
 \label{eq:eic-residual-reflection}
\end{equation}
For the forward direction, substitute $x=Sa$, $y=b$ into
\eqref{eq:eic-residual-invariant}.  If $d(Sa,b)=1$, the sign is $1$
and the selector returns its second branch, giving
\[
 b+1=Sa.
\]
The defining equation for addition gives $b+1=S(b+0)=Sb$.  Thus
$Sb=Sa$, and applying $P$ gives $b=a$, hence $a=b$.
Conversely, $a=b$ and \eqref{eq:eic-diagonal} give
$d(Sa,b)=d(Sa,a)=1$.

Define the equality characteristic by the fixed description
\begin{equation}
 e(a,b)=\operatorname{sg}\bigl(K(d(Sa,b))\bigr).
 \label{eq:eic-equality-characteristic}
\end{equation}
The sign, countdown, and residual reflection lemmas give
\begin{align}
 e(a,b)=0&\quad\longleftrightarrow\quad a=b,
       \label{eq:eic-equality-reflection}\\
 e(a,b)=0&\ \lor\ e(a,b)=1.
       \label{eq:eic-equality-range}
\end{align}
In detail, $e(a,b)=0$ is equivalent to $K(d(Sa,b))=0$ by
\eqref{eq:eic-sign-zero}, which is equivalent to $d(Sa,b)=1$ by
\eqref{eq:eic-countdown-zero}, and then to $a=b$ by
\eqref{eq:eic-residual-reflection}.  The range assertion is
\eqref{eq:eic-sign-range} applied to $K(d(Sa,b))$.

The reflection property concerns the logic's actual equality throughout
the domain.  It was not assumed as part of a numerical comparison symbol,
and its proof did not restrict $a$ or $b$ to standard numerals.
It also yields decidable equality.  In the $e(a,b)=0$ branch,
\eqref{eq:eic-equality-reflection} gives $a=b$.  In the $e(a,b)=1$
branch, an assumption $a=b$ would give $e(a,b)=0$ and therefore $0=1$,
contradicting separation.  Hence
\begin{equation}
 \forall a\,\forall b\,\bigl(a=b\lor(a=b\to\bot)\bigr)
 \label{eq:eic-equality-decidable}
\end{equation}
is a theorem, not an extra hypothesis.

\subsection{Compilation of all quantifier-free formulas}

For each fixed quantifier-free formula $A(\vec p)$, construct a term
$b_A(\vec p)$ by external recursion on its finite syntax.  Take equality,
$\bot$, $\land$, $\lor$, and $\to$ as the primitive formula constructors;
negation is implication into $\bot$.  Set
\begin{align}
 b_{s=t}&=e(s,t), & b_{\bot}&=1,\label{eq:eic-boolean-atoms}\\
 b_{A\land D}&=C(b_A,b_D,1),\label{eq:eic-boolean-and}\\
 b_{A\lor D}&=C(b_A,0,b_D),\label{eq:eic-boolean-or}\\
 b_{A\to D}&=C(b_A,b_D,0).\label{eq:eic-boolean-implies}
\end{align}
Every term here is a finite composition of fixed primitive-recursive
descriptions; if desired, projections and composition give it a single
description symbol.  No internal quantification over formulas or programs
is involved.

An external induction on the syntax of $A$ proves, within the arithmetic
theory,
\begin{equation}
 \bigl(b_A=0\lor b_A=1\bigr)
 \quad\text{and}\quad
 \bigl(b_A=0\longleftrightarrow A\bigr).
 \label{eq:eic-boolean-correctness}
\end{equation}
Here are the full connective checks.  In every case $b_A=1$ implies
$\neg A$: if $A$ held, the already established implication
$A\to b_A=0$ would contradict separation.

\paragraph{Atomic formulas and falsity.}
For an equation, both claims are
\eqref{eq:eic-equality-reflection}--\eqref{eq:eic-equality-range}.
For falsity the range assertion is $1=0\lor1=1$, with its right disjunct
proved by reflexivity.  Also $1=0\to\bot$ follows from separation and
symmetry, while $\bot\to1=0$ follows by explosion.  Thus $1=0$ is
equivalent to $\bot$, as required.

\paragraph{Conjunction.}
Split $b_A=0\lor b_A=1$.  In the zero branch, $A$ holds and
$b_{A\land D}=b_D$.  Its range follows from the range of $b_D$,
and its being zero is equivalent to $D$, hence to $A\land D$ under
the known $A$.  In the one branch, $\neg A$ holds and
$b_{A\land D}=1$.  Its range is immediate.  If it were zero,
separation would give $\bot$, hence $A\land D$; conversely an
assumption $A\land D$ contradicts $\neg A$, and explosion gives
$b_{A\land D}=0$.

\paragraph{Disjunction.}
In the branch $b_A=0$, the formula $A$ holds and
$b_{A\lor D}=0$.  The range is immediate, and both $b_{A\lor D}=0$
and $A\lor D$ hold, proving their equivalence.  In the branch
$b_A=1$, we have $\neg A$ and $b_{A\lor D}=b_D$.  Its range follows
from the range of $b_D$.  If its value is zero, then $D$ and hence
$A\lor D$.  Conversely, from $A\lor D$, the $A$ alternative
contradicts $\neg A$ and yields $b_D=0$ by explosion, while the
$D$ alternative yields $b_D=0$ directly.  This proves the required
equivalence in both alternatives.

\paragraph{Implication.}
In the branch $b_A=0$, the formula $A$ holds and
$b_{A\to D}=b_D$.  Its range follows from the range of $b_D$.
If its value is zero, then $D$, which gives $A\to D$.  Conversely,
$A\to D$ and the known $A$ give $D$, hence $b_D=0$.
In the branch $b_A=1$, we have $\neg A$ and
$b_{A\to D}=0$.  Its range is immediate.  To prove $A\to D$,
assume $A$; this contradicts $\neg A$, and explosion gives $D$.
Thus both $b_{A\to D}=0$ and $A\to D$ hold in this branch, and
their equivalence follows.  This argument uses the proved two-value
range of the finite characteristic.  It does not use excluded middle
for arbitrary formulas or change the surrounding intuitionistic logic.

This completes the external syntactic induction proving
\eqref{eq:eic-boolean-correctness} for every quantifier-free formula.

\subsection{Recovery of induction and the full theory}

Let $A(n,\vec p)$ be any fixed quantifier-free formula.  Suppose an
arbitrary retained context $\Gamma$ proves its base and its induction
step:
\[
 \Gamma\vdash A(0,\vec p),
 \qquad
 \Gamma,A(z,\vec p)\vdash A(Sz,\vec p),
\]
with the usual freshness condition on $z$.  By
\eqref{eq:eic-boolean-correctness}, the base gives
$\Gamma\vdash b_A(0,\vec p)=0$.  Under the additional assumption
$b_A(z,\vec p)=0$, the same equivalence gives $A(z,\vec p)$,
the given step gives $A(Sz,\vec p)$, and the equivalence gives
$b_A(Sz,\vec p)=0$.  Apply literal-equation induction to
$b_A(n,\vec p)=0$ and then apply its reflection implication.
The result is $\Gamma\vdash A(t,\vec p)$ for every target term $t$.
Thus the theory proves the full contextual quantifier-free induction
rule, including in arbitrary quantified contexts.

\paragraph{Completion theorem.}
In the common language of fixed primitive-recursive descriptions, the
theory
\[
 \mathsf{PRE}+\mathsf{Cov}+\mathsf{Sep}+\mathsf{EI}
\]
has exactly the same sequents as ordinary first-order intuitionistic
primitive recursive arithmetic with all quantifier-free induction
instances.  The forward simulation follows from the recovered
quantifier-free induction, successor injectivity, and successor separation;
all primitive-recursive defining equations and all logical rules are
shared.  For the reverse simulation, literal-equation induction is an
instance of quantifier-free induction, and successor separation gives
$\mathsf{Sep}$.  Quantifier-free induction proves $\mathsf{Cov}$:
at zero choose its left disjunct, and at a successor use
$Sx=S(P(Sx))$, which follows from $P(Sx)=x$, for its right disjunct.

Both interpretations therefore use the full domain, defined by $x=x$,
and preserve every function symbol, equality, logical connective, and
quantifier.  Both composites have the identity isomorphism, with graph
$x=y$.  The theories are consequently identity-domain bi-interpretable
and equiconsistent.  This completion theorem is unconditional once the
literal-equation induction rule has been obtained; it has no dependency
on a classical successor-uniqueness reduction.

% END include

% BEGIN appendices/primitive_recursion.tex
\section{Pairing, primitive recursion, and internal numerical composites}
\label{sec:iterator-coding}
This appendix is adapted with attribution from the progress-closure
arithmetic development~\cite{progress}. Its prerequisites have now been
proved: scalar predecessor, arithmetic reflection, and full quantifier-free
induction. Pairing is a derived proof tool; it is not part of any checker.

\subsection{Pairing from a bounded square-root search}

Once QFI is available, a familiar square-based injection gives a
particularly direct compiler. Write $a\le b$ for $d(a,b)=0$.
The usual linear-order and semiring identities follow by fixed
quantifier-free inductions from the successor, addition, and monus
equations. Multiplication itself is the scalar term
\[
                        a\cdot b=I(b,0;u.u+a).
\]
Only primitive-recursive arithmetic is involved in the following
bounded root search; there is no minimization over an unbounded domain.

For a fixed parameter $n$, put
\[
 U_n(v)=C(d((Sv)^2,n),Sv,v),\qquad
 R_n(i)=I(i,0;v.U_n(v)),\qquad q(n)=R_n(n).
\]
The update increments $v$ exactly when $(v+1)^2\le n$, and
otherwise retains it. QFI in $i$, with $n$ fixed, proves
\[
 R_n(i)\le i,\qquad R_n(i)^2\le n,\qquad
 R_n(i)=i\ \vee\ n<(S(R_n(i)))^2.              \tag{Root invariant}
\]
The base has $R_n(0)=0$. If the update increments, its test
contradicts the strict-upper-bound alternative in the third invariant.
Thus the old state equals $i$, and the new state equals $Si$;
its square bound is the increment test. If the update does not increment, the
strict upper square bound holds and the state remains fixed.
An already established strict upper square bound forces every
subsequent update to remain fixed. These are exactly the cases
needed for the three displayed invariants.

At $i=n$, either the strict upper square bound already holds or
$R_n(n)=n$, in which case $n<(n+1)^2$ supplies it. Thus
\[
                        q(n)^2\le n<(q(n)+1)^2.           \tag{Root}
\]
Monotonicity of squares shows uniqueness of a number satisfying
Root. All formulas just used are quantifier free in the scalar
language, with fixed $n$ during the induction.

Now define
\[
 \pi(a,b)=(a+b)^2+a,\qquad
 L(n)=d(n,q(n)^2),\qquad H(n)=d(q(n),L(n)).
\]
For $k=a+b$, one has $k^2\le k^2+a<(k+1)^2$ since $a\le k$.
Root uniqueness gives $q(\pi(a,b))=a+b$, and therefore
\[
                    L(\pi(a,b))=a,\qquad H(\pi(a,b))=b.
\]
Surjectivity of $\pi$ is unnecessary. The original numerical domain
is never replaced by the range of this coding; pairing is used only
to store the finite state of a fixed primitive recursion. This
standard square pairing is a convenient reconstruction here, not
a claimed new arithmetic basis.

\subsection{Compiling each primitive recursion}

Suppose the fixed target description is
\[
 f(0,\vec p)=g(\vec p),\qquad
 f(Sn,\vec p)=h(n,f(n,\vec p),\vec p).
\]
Assume by external structural induction that the earlier $g,h$
already have scalar translations $g^*,h^*$. Define
\[
 W(u,\vec p)=\pi(SL(u),h^*(L(u),H(u),\vec p)),
\]
\[
 V(n,\vec p)=I(n,\pi(0,g^*(\vec p));u.W(u,\vec p)),
 \qquad f^*(n,\vec p)=H(V(n,\vec p)).
\]
The update has only its state variable bound, with every $p_i$
fixed. Equation induction proves $L(V(n,\vec p))=n$.
The two projection equations and scalar computation then prove
both defining equations for $f^*$ internally. Initial functions
and composition translate directly. This completes the translation
of every fixed finite PR description without adding a universal
evaluator.

\subsection{Why both numerical round trips are internal}

Two fixed scalar iterators with the same seed and provably equal
updates are internally equal at all counters: use literal-equation
induction comparing their values. Their base agrees, and their
successor equations plus congruence prove the step. The same
uniqueness argument works for ordinary primitive recursion.

For an ordinary PR description, induct externally on its finite
construction. The double translation satisfies its original
defining equations, so internal equation induction identifies it
with the original operation. For a native scalar description,
earlier translated updates agree by the external induction
hypothesis; scalar uniqueness identifies the double-translated
iterator. These proofs hold at arbitrary internal counters.

Thus the domain relation $x=y$ is provably a total bijection
preserving every fixed numerical operation in both composites.
The formula constructions and round trips in
Section~\ref{sec:choice-introduction} give the corresponding formula
composites. Both interpretations preserve falsity, taking native
$0=1$ to an ordinary contradiction and ordinary $\bot$ to native
$0=1$. Finite proof translations therefore establish equiconsistency.

% END include

% BEGIN appendices/inherited_ordinary_bridge.tex
\section{Arithmetic strength and both ordinary interpretation composites}
\label{sec:ordinary-bridge}

We spell out the arithmetic bridge to avoid identifying standard numerical
agreement with bi-interpretability. First elaborate every fixed iterator
into an ordinary function symbol. At stage zero the symbols are $0,S$.
Given a fixed earlier-stage term $h(u,\vec p)$, introduce
$i_h(n,a,\vec p)$ with the two equations
\[
i_h(0,a,\vec p)=a,\qquad
i_h(Sn,a,\vec p)=h(i_h(n,a,\vec p),\vec p).
\]
Only finite descriptions occur; the domain contains no code for a
universal evaluator. Formal parameters are distinct when a description
is declared, while actual arguments can coincide.

Let $T_I$ be ordinary intuitionistic first-order logic with actual
equality and genuine falsity, these staged equations, contextual
induction for literal equations, and
\[
\forall x\,(x=0\vee\exists y\,x=Sy),
\qquad (0=S0)\to\bot.
\]
The second sentence is the falsity bridge. It must be stated: positive
numerical equations by themselves admit a singleton structure. Native
explosion from $0=S0$ and a translation of ordinary falsity to that
equation do not authorize silently dropping the bridge from $T_I$.

The native-to-ordinary translation is
\[
(s=t)^\circ=(s=t),\qquad
(\mathsf D_{x,y}B)^\circ=
\exists x\bigl(B^\circ(x,0)\to\forall y B^\circ(x,y)\bigr).
\]
A clause $\langle A;y.B\rangle$ becomes $A^\circ\to\forall yB^\circ$.
For the auxiliary closure used by the compact checkers,
$(\mathsf Q_yB)^\circ=B^\circ(0)\to\forall yB^\circ$.
The interpretation domain is all elements, defined by $x=x$;
equality is actual equality.

Case DI is sound because numerical coverage combines its retained base
with its successor proof. Progress DI is sound by literal-equation
induction under a temporarily assumed base. Use, DE and transport are
ordinary logical inferences under the translations. Conversely the
derived logical operations, equation induction, cases and explosion in
Section~\ref{sec:choice-introduction} recover the rules and axioms of
$T_I$. An auxiliary closure is interderivable with
$\mathsf D_{v,y}B(y)$ for unused $v$: in one direction use case DI with
the closure and its base; in the other eliminate the paired formula
and return its supplied closure. Thus formulas and context clauses have
the stated meanings in both presentations.

\subsection{The arithmetic dependency order}

The scalar update has no recursion-counter argument. We do not conceal
one in an encoding. The inherited internal proofs, reproduced in the
appendices, establish the required arithmetic in this order:
\begin{enumerate}
\item Using equation induction and numerical cases, construct a fixed
scalar predecessor term and prove $P0=0$ and $P(Sn)=n$.
Existential coverage then gives $x=0\vee x=S(Px)$: in the successor
case $x=Sy$, apply $P$ to obtain $Px=y$ and substitute.
\item Construct truncated subtraction and a numerical characteristic
of actual equality. Prove its reflection on the entire original domain.
\item Compile each fixed quantifier-free formula into its characteristic
equation and recover every contextual quantifier-free induction instance.
\item With that induction available, construct pairing and projections
by a bounded square-root search and compile every fixed general
primitive-recursive description into a scalar iterator term.
\end{enumerate}
No pairing or faithful equality detector is assumed to obtain
predecessor. No appeal to standard-model correctness replaces an
internal equality. Every induction fixes its passive parameters;
diagonal substitution occurs only into a previously proved identity.

For a target recursion
\[
f(0,\vec p)=g(\vec p),\qquad
f(Sn,\vec p)=h(n,f(n,\vec p),\vec p),
\]
let $\pi,L,H$ be the proved pairing and projections. The compiler uses
\[
V(n,\vec p)=I\bigl(n,\pi(0,g^*(\vec p));
u.\pi(SL(u),h^*(L(u),H(u),\vec p))\bigr),
\qquad f^*(n,\vec p)=H(V(n,\vec p)).
\]
Equation induction proves $L(V(n,\vec p))=n$; the two projection
identities then prove the target defining equations. The construction
is external recursion on each finite PR description, not an internal
quantification over arbitrary functions.

\begin{theorem}[Whole-domain bi-interpretability and equiconsistency]
The ordinary realization $T_I$ and full first-order iPRA are
bi-interpretable. In both theories $x=y$ defines the isomorphism from
the original domain to its double interpretation. The native Choice-DI
calculus is therefore a faithful proof presentation of a theory
bi-interpretable and equiconsistent with iPRA.
\end{theorem}
\begin{proof}
Interpret an iterator in iPRA by primitive recursion that ignores its
counter argument in the update. iPRA proves its equations, literal
equation induction, coverage and separation. In the reverse direction
use the scalar compiler. The predecessor and equality-reflection proofs
supply the successor axioms and every quantifier-free induction instance.
All ordinary logical rules are preserved.

For two iterators with the same seed and provably equal updates,
equation induction on the equality of their values proves uniqueness at
every counter, with all parameters fixed. The corresponding argument
applies to ordinary primitive recursion. Induct externally on a finite
description: its double translation satisfies its original defining
equations, so internal uniqueness identifies the two functions. This
also proves coherence when substitution changes the staged syntactic
description of an update.

The domain and equality are unchanged; hence the identity graph is
internally total, bijective and preserves every interpreted operation.
The displayed ordinary meanings of the logical macros and the paired
binder give both formula composites by external induction on finite
formula syntax. Clauses are similarly equivalent to their conditional
universal meanings. The explicit falsity bridge preserves contradiction
in both directions. Effective translations of finite proofs give
equiconsistency, without a claim that either theory proves its own
consistency.
\end{proof}

% END include


\clearpage
\begin{thebibliography}{99}\raggedright
\bibitem{snapshot}
Vladimir Reshetnikov, repository maintainer. \emph{ProveIt: current draft
collection}, commit \file{4a04139595eadc3e8818f1df2a94838ed4851d83}.
\url{https://github.com/VladimirReshetnikov/ProveIt/tree/4a04139595eadc3e8818f1df2a94838ed4851d83/docs/drafts}.

\bibitem{generated}
\emph{Numerical axioms as local assumptions},
\file{ipra_generated_assumptions_2026-11-10.zip} in~\cite{snapshot}.
Source of the current leading generator, zero-witness discharge,
scope decoder, test producer, and reproduced bridge.

\bibitem{named}
\emph{Naming values before specialization},
\file{ipra_named_specialization_2026-11-10.zip} in~\cite{snapshot}.
Source of the named-specialization comparison and the question about
simultaneous restrictions.

\bibitem{endpoint}
\emph{Endpoint certificates},
\file{ipra_endpoint_certificates_2026-11-09.zip} in~\cite{snapshot}.
Source of the shared whole-program measurement convention.

\bibitem{erasure}
Inherited scalar predecessor construction, preserved as
\file{reference/scalar_predecessor.tex} in the generated-assumptions
archive, with its earlier scalar-development attribution.

\bibitem{forward}
Inherited actual-equality reflection construction, preserved as
\file{reference/equality_reflection.tex} in the generated-assumptions
archive, with attribution to the forward-left development.

\bibitem{progress}
Inherited pairing and primitive-recursion compiler, preserved as
\file{reference/primitive_recursion.tex} in the generated-assumptions
archive, with attribution to the progress-closure development.

\bibitem{salehi}
Saeed Salehi. \emph{Decidable Formulas of Intuitionistic Primitive
Recursive Arithmetic}. Reports on Mathematical Logic \textbf{36}
(2002), 55--61.
\url{https://rml.tcs.uj.edu.pl/rml-36/36-SALEHI.pdf}.
\end{thebibliography}
\end{document}
